\documentclass[aspectratio=169,10pt]{beamer}
% notes.tex defines \NOTESONLY and reads this file: then only the speaker notes are printed
\ifdefined\NOTESONLY
  \setbeameroption{show only notes}
\fi

\usepackage{lmodern}
\usepackage[utf8]{inputenc}
\usepackage[T1]{fontenc}
\usepackage[polish,english]{babel}
\usepackage{amsmath,amssymb}
\usepackage{booktabs}
\usepackage{colortbl}
\usepackage{xcolor}
\usepackage{fontawesome5}
\usepackage{listings}
\usepackage{etoolbox}
\usepackage{tikz}
\usetikzlibrary{arrows.meta, calc, positioning, fit, decorations.pathreplacing, decorations.pathmorphing}
\usepackage{pgfplots}
% layers for the code windows drawn in TikZ: the window, then the line that runs now, then the text
\pgfdeclarelayer{codebg}\pgfdeclarelayer{codehl}\pgfsetlayers{codebg,codehl,main}
\pgfplotsset{compat=1.18}

\usetheme{Madrid}
\useinnertheme{rounded}
\useoutertheme{miniframes}
\setbeamertemplate{navigation symbols}{}
\setbeamertemplate{itemize items}[circle]
\setbeamertemplate{enumerate items}[default]
\setbeamertemplate{blocks}[rounded][shadow=true]
% one dot per idea, not per frame: \samedot before a frame gives it the dot of the frame before
% (it undoes the step of the dot counter that the next frame makes)
\newcommand{\samedot}{\addtocounter{subsectionslide}{-1}}
% the default dot step of the circle dots (0.1 cm dots); offset -.04 keeps the old header height (-.01 made 2 slides overfull)
\setbeamersize{mini frame size=.14cm,mini frame offset=-.04cm}

% ======================================================================
% Colour rules for presentations (Beamer + TikZ), from Agents/colors
% ======================================================================

% ---------- template ----------
\definecolor{DeepBlue}{HTML}{283593}      % title bar, headings, links, control-flow arrows
\definecolor{AccentOrange}{HTML}{EB5A14}  % annotations: step numbers, "look here", code highlight, conclusions
\definecolor{LightGray}{HTML}{F5F6FA}     % background of information blocks
\setbeamercolor{palette primary}{bg=DeepBlue,fg=white}
\setbeamercolor{palette secondary}{bg=DeepBlue!80!black,fg=white}
\setbeamercolor{palette tertiary}{bg=DeepBlue!90!black,fg=white}
\setbeamercolor{title}{bg=DeepBlue,fg=white}
\setbeamercolor{frametitle}{bg=DeepBlue,fg=white}
\setbeamercolor{structure}{fg=DeepBlue}
\setbeamercolor{item}{fg=AccentOrange}
\setbeamercolor{alerted text}{fg=AccentOrange}
\setbeamercolor{block title}{bg=DeepBlue,fg=white}
\setbeamercolor{block body}{bg=LightGray,fg=black}
\setbeamercolor{block title example}{bg=AccentOrange,fg=white}
\setbeamercolor{block body example}{bg=AccentOrange!10,fg=black}
\setbeamercolor{enumerate item}{fg=DeepBlue}
\setbeamerfont{enumerate item}{series=\bfseries}

% ---------- actors and categories: fills only (S1 to S6) ----------
\definecolor{sIndigo}{HTML}{5C6BC0}   % S1
\definecolor{sPink}{HTML}{EC407A}     % S2
\definecolor{sViolet}{HTML}{7E57C2}   % S3
\definecolor{sAmber}{HTML}{FFB300}    % S4: black text on it
\definecolor{sTeal}{HTML}{26A69A}     % S5
\definecolor{sBrown}{HTML}{8D6E63}    % S6
% text variants that stay readable on white
\definecolor{sPinkText}{HTML}{C2185B}
\definecolor{sAmberText}{HTML}{A06A00}
\definecolor{sTealText}{HTML}{00796B}

% ---------- actions: frames only; results: signs only ----------
\definecolor{aExec}{HTML}{43A047}     % runs now
\definecolor{aRead}{HTML}{1E88E5}     % is read
\definecolor{aWrite}{HTML}{E53935}    % is written

% ---------- neutral (Blue Grey) ----------
\definecolor{nInk}{HTML}{37474F}      % text in figures, waveforms, wires, data-flow arrows
\definecolor{nLine}{HTML}{90A4AE}     % borders, axes
\definecolor{nLight}{HTML}{CFD8DC}    % inactive lines, grids, empty places (dashed)
\definecolor{nFill}{HTML}{ECEFF1}     % fill of boxes and used cells
\definecolor{nIdle}{HTML}{E0E0E0}     % time lost (idle, busy waiting, bubbles)

% ---------- TikZ styles ----------
\tikzset{
  aexec/.style={draw=aExec, line width=1.6pt},
  aexecm/.style={aexec, label={[text=aExec, inner sep=1pt, font=\tiny]left:$\blacktriangleright$}},
  aread/.style={draw=aRead, line width=1.6pt, dash pattern=on 3pt off 2pt},
  awrite/.style={draw=aWrite, line width=2.4pt},
  hw/.style={draw=nLine, thick, fill=nFill, text=nInk, rounded corners=3pt, align=center, inner sep=4pt},
  hwmain/.style={draw=nInk, thick, fill=nInk, text=white, rounded corners=3pt, align=center, inner sep=4pt},
  cell/.style={draw=nLine, thick, fill=white, text=nInk, font=\ttfamily\scriptsize},
  cellused/.style={cell, fill=nFill},
  cellfree/.style={cell, dashed, draw=nLight},
  container/.style={draw=nLine, thick, dashed, rounded corners=5pt},
  actor/.style={fill=#1, draw=#1, text=white, rounded corners=3pt, align=center, font=\scriptsize\bfseries, inner sep=4pt},
  actorlight/.style={fill=#1!15, draw=#1, thick, text=nInk, rounded corners=3pt, align=center, inner sep=4pt},
  suspended/.style={fill=#1!30, draw=#1, dashed, line width=0.9pt, text=nInk, rounded corners=3pt, align=center, inner sep=4pt},
  ready/.style={fill=white, draw=#1, line width=0.8pt, text=nInk, rounded corners=3pt, align=center, inner sep=4pt},
  idle/.style={fill=nIdle, draw=none, text=nInk, rounded corners=3pt, align=center, inner sep=4pt},
  flow/.style={-{Stealth}, very thick, DeepBlue},
  data/.style={-{Stealth}, very thick, nInk},
  holds/.style={-{Stealth}, very thick, nInk},
  waits/.style={-{Stealth}, very thick, nLine, dashed},
  annot/.style={draw=AccentOrange, line width=1.4pt, rounded corners=3pt},
}
\tikzset{actoramber/.style={actor=sAmber, text=black}}

% ---------- signs and small helpers ----------
\newcommand{\ok}{\textcolor{aExec}{\faCheck}}
\newcommand{\nok}{\textcolor{aWrite}{\faTimes}}
\newcommand{\badge}[1]{\tikz[baseline=(b.base)]\node[circle,fill=AccentOrange,text=white,
  inner sep=1.2pt,font=\bfseries\tiny] (b) {#1};}
\newcommand{\bigbadge}[1]{\tikz[baseline=(b.base)]\node[circle,fill=AccentOrange,text=white,
  inner sep=2pt,font=\bfseries\small] (b) {#1};}
% ---------- text boxes ----------
\newcommand{\defbox}[1]{%
  \begin{beamercolorbox}[rounded=true,sep=0.1cm,wd=\linewidth]{block body example}
    \footnotesize #1
  \end{beamercolorbox}\vspace{0.06cm}}
% a large definition in the middle of the slide
\newcommand{\bigdef}[1]{%
  \begin{beamercolorbox}[rounded=true,sep=0.25cm,wd=\linewidth]{block body example}
    \centering\normalsize #1
  \end{beamercolorbox}\vspace{0.06cm}}
\newcommand{\stepbox}[2]{%
  \begin{beamercolorbox}[rounded=true,sep=0.09cm,wd=\linewidth]{block body}
    \footnotesize\textbf{#1}\\ #2
  \end{beamercolorbox}\vspace{0.06cm}}
\newcommand{\keymsg}[1]{%
  \begin{beamercolorbox}[rounded=true,sep=0.1cm,wd=\textwidth]{block body example}
    \centering\footnotesize #1
  \end{beamercolorbox}}
% a grey box without a title
\newcommand{\plainbox}[1]{%
  \begin{beamercolorbox}[rounded=true,sep=0.12cm,wd=\linewidth]{block body}
    \centering\footnotesize #1
  \end{beamercolorbox}\vspace{0.06cm}}
\setbeamercolor{block body problem}{bg=AccentOrange!10,fg=black}
\newcommand{\probbox}[2]{%
  \begin{beamercolorbox}[rounded=true,sep=0.09cm,wd=\linewidth]{block body problem}
    \footnotesize\textbf{\textcolor{AccentOrange}{\faExclamationTriangle}\ #1}\\ #2
  \end{beamercolorbox}\vspace{0.06cm}}

% ---------- code: dark window, grey dots, syntax colours ----------
\definecolor{VSDarkBg}{RGB}{25,25,25}
\definecolor{VSText}{RGB}{245,245,245}
\definecolor{VSBlue}{RGB}{80,180,255}
\definecolor{VSPurple}{RGB}{210,110,210}
\definecolor{VSString}{RGB}{255,140,80}
\definecolor{VSComment}{RGB}{80,220,80}
\setbeamercolor{vscodelisting}{bg=VSDarkBg}
\lstdefinestyle{code}{basicstyle=\color{VSText}\ttfamily\scriptsize, keywordstyle=\color{VSBlue}\bfseries,
  stringstyle=\color{VSString}, commentstyle=\color{VSComment}\itshape, emphstyle=\color{VSPurple},
  numbers=none, columns=fixed, basewidth=0.52em, keepspaces=true, showstringspaces=false, xleftmargin=6pt}
\lstdefinestyle{py}{style=code, language=Python, emph={print,chr,zip,range,len,bytes,ord,hex}}
\BeforeBeginEnvironment{lstlisting}{%
  \begin{beamercolorbox}[rounded=true,sep=2mm,wd=\linewidth]{vscodelisting}%
  \begin{tikzpicture}
    \fill[nLine] (0,0) circle (3pt); \fill[nLine] (0.35,0) circle (3pt); \fill[nLine] (0.70,0) circle (3pt);
  \end{tikzpicture}%
  \vspace{-0.2cm}}
\AfterEndEnvironment{lstlisting}{\end{beamercolorbox}}
\newcommand{\hlcode}[1]{{\setlength{\fboxsep}{1pt}\colorbox{AccentOrange}{\textcolor{white}{\bfseries #1}}}}
\newcommand{\eol}{\tikz[remember picture, overlay]\coordinate (eol) at (0,0.08);}
\newcommand{\codenote}[3][0pt]{%
  \begin{tikzpicture}[remember picture, overlay, >=Stealth]
    \node[draw=AccentOrange, very thick, fill=AccentOrange!6, rounded corners=4pt, text width=0.33\paperwidth,
      align=left, inner sep=6pt, font=\footnotesize, anchor=west] (cal)
      at ([yshift=#1]eol -| {$(current page.west)!0.6!(current page.east)$}) {\textbf{\textcolor{AccentOrange}{#2}}\\[3pt]#3};
    \draw[->, very thick, AccentOrange] (eol -| cal.west) -- ([xshift=4pt]eol);
  \end{tikzpicture}}
% program output: a light terminal line
\newcommand{\terminal}[1]{%
  \begin{tikzpicture}
    \node[draw=nLine, thick, fill=nFill, rounded corners=3pt, inner sep=5pt, font=\ttfamily\small, text=nInk]
      {\textcolor{nLine}{\$}\ #1};
  \end{tikzpicture}}
% the Python prompt: \pyprompt{pow(5, 2, 23)}{2} prints >>> pow(5, 2, 23) and the result line below it
\newcommand{\pyprompt}[2]{%
  \begin{tikzpicture}[baseline={([yshift=-0.5ex]t.center)}]
    \node[draw=nLine, thick, fill=nFill, rounded corners=3pt, inner sep=5pt, font=\ttfamily\small, text=nInk, align=left] (t)
      {\textcolor{nLine}{>{}>{}>}\ #1\\#2};
  \end{tikzpicture}}

% ======================================================================
% Course macros
% ======================================================================
% the starter files on GitHub: the word "GitHub" is the link.
%   \ghlink{Task4} opens the folder of Task 4, \ghlink{} opens the folder of this lecture
\newcommand{\ghbase}{https://github.com/KsaweryM/Cryptography-Course/tree/main/Lecture04-Diffie-Hellman}
\newcommand{\ghlink}[1]{\href{\ghbase/#1}{\textcolor{DeepBlue}{\textbf{GitHub}}}}
% the task slides: the steps are in the README on GitHub, the slide only opens the task
%   \taskgoal{one line}   \taskready{the ready modules}   \opentask{Task2}{about 6 minutes}
%   a task on paper: the question on the slide, then \onpaper{about 5 minutes}
\newcommand{\taskgoal}[1]{{\centering\large\textcolor{AccentOrange}{\faBullseye}\ \textbf{Goal:} #1\par}\vspace{0.2cm}}
\newcommand{\taskready}[1]{{\centering\normalsize\textcolor{nInk}{\faBox}\ \textbf{Ready:} #1\par}\vspace{0.2cm}}
\newcommand{\opentask}[2]{%
  \begin{beamercolorbox}[rounded=true,sep=0.35cm,wd=\linewidth]{block body example}
    \centering
    {\LARGE\faGithub\ \textbf{Open the task:}\ \href{\ghbase/#1}{\textcolor{DeepBlue}{\textbf{\texttt{#1}}}}}\\[8pt]
    \footnotesize\href{\ghbase/#1}{\textcolor{DeepBlue}{\texttt{github.com/KsaweryM/Cryptography-Course/tree/main/Lecture04-Diffie-Hellman/#1}}}\\[6pt]
    \small on \ghlink{#1}: the steps, the starter files and how to run them\ifx&#2&\else\qquad\faClock[regular]\ #2\fi
  \end{beamercolorbox}}
\newcommand{\onpaper}[1]{%
  \begin{beamercolorbox}[rounded=true,sep=0.3cm,wd=\linewidth]{block body example}
    \centering
    {\LARGE\faPencil*\ \textbf{On paper}}\ifx&#1&\else\qquad\small\faClock[regular]\ #1\fi
  \end{beamercolorbox}}

% a page of the simulator website, printed in full so that students can type it:
%   \simlink{diffie-hellman/\#1} prints and links ksawerym.github.io/Cryptography-simulators/diffie-hellman/#1
\newcommand{\simbase}{ksawerym.github.io/Cryptography-simulators}
\newcommand{\simlink}[1]{\href{https://\simbase/#1}{\textcolor{DeepBlue}{\texttt{\simbase/#1}}}}
% "For the curious" line under a frame: \curious{diffie-hellman/\#1}
\newcommand{\curious}[1]{{\footnotesize\textcolor{gray}{\faSearch\ For the curious:}\ \simlink{#1}}}

% speaker notes (Polish): printed only in the notes PDF (make notes)
\setbeamertemplate{note page}[default]
\setbeamerfont{note page}{size=\normalsize}
\newcommand{\say}[1]{\note{\selectlanguage{polish}#1}}

\usetikzlibrary{shadows}
% long numbers that may break anywhere
\usepackage{seqsplit}

% ======================================================================
% Colour mapping for this lecture (colour rules, section 6)
%   S1 sIndigo : Alice
%   S2 sPink   : the attackers (Eve listens, Mallory changes)
%   S3 sViolet : the message M
%   S4 sAmber  : secret values: a, b, x, k, m and the shared secret
%   S5 sTeal   : Bob
%   S6 sBrown  : ciphertext (C1, C2)
%   public values (p, g, A, B, y) are neutral: grey boxes
%   paint colours in the mixing story are pictures of paint, not roles
% ======================================================================
\definecolor{paintY}{HTML}{FDD835}
\definecolor{paintR}{HTML}{E53935}
\definecolor{paintB}{HTML}{1E88E5}
\definecolor{paintO}{HTML}{FB8C00}
\definecolor{paintG}{HTML}{7CB342}
\definecolor{paintM}{HTML}{6D4C41}
\tikzset{
  pub/.style={draw=nLine, thick, fill=nFill, text=nInk, rounded corners=3pt, inner sep=4pt, align=center},
  sec/.style={draw=sAmber, thick, fill=sAmber!25, text=nInk, rounded corners=3pt, inner sep=4pt, align=center},
  msgb/.style={draw=sViolet, thick, fill=sViolet!15, text=nInk, rounded corners=3pt, inner sep=4pt, align=center},
  ctx/.style={draw=sBrown, thick, fill=sBrown!15, text=nInk, rounded corners=3pt, inner sep=4pt, align=center},
  net/.style={very thick, nInk, decorate, decoration={snake, amplitude=1.4pt, segment length=8pt}},
  netarrow/.style={-{Stealth}, very thick, nInk, decorate, decoration={snake, amplitude=1.4pt, segment length=8pt, post length=5pt}},
  capt/.style={font=\scriptsize, text=nInk, align=center},
  rowlabel/.style={font=\scriptsize, text=nInk, anchor=east},
  rowhl/.style={fill=AccentOrange!10, draw=AccentOrange, thick, rounded corners=3pt},
}
% a person: \who{x}{y}{fill}{text colour}{name}{icon}; the node is called by the name
\newcommand{\who}[6]{%
  \node[circle, fill=#3, minimum size=1.1cm, inner sep=0pt, text=white, font=\Large] (#5) at (#1,#2) {#6};
  \node[font=\small\bfseries, text=#4, below=1pt of #5] (#5name) {#5};}
\newcommand{\alice}[2]{\who{#1}{#2}{sIndigo}{sIndigo}{Alice}{\faUser}}
\newcommand{\bob}[2]{\who{#1}{#2}{sTeal}{sTealText}{Bob}{\faUser}}
\newcommand{\eve}[2]{\who{#1}{#2}{sPink}{sPinkText}{Eve}{\faUserSecret}}
\newcommand{\mallory}[2]{\who{#1}{#2}{sPink}{sPinkText}{Mallory}{\faUserNinja}}
% an envelope: \envpic{x}{y}{fill}{border}
\newcommand{\envpic}[4]{%
  \begin{scope}[shift={(#1,#2)}]
    \draw[#4, thick, fill=#3, rounded corners=1pt] (-0.8,-0.5) rectangle (0.8,0.5);
    \draw[#4, thick, line join=round] (-0.8,0.5) -- (0,-0.05) -- (0.8,0.5);
  \end{scope}}
% a step list: \stepitem{number}{current step}{text}
\newcommand{\stepitem}[3]{%
  \ifnum#1=#2\relax
    \textcolor{AccentOrange}{$\blacktriangleright$}\ \textbf{#1. #3}%
  \else\ifnum#1<#2\relax
    \phantom{$\blacktriangleright$}\ \textcolor{nInk}{#1. #3}%
  \else
    \phantom{$\blacktriangleright$}\ \textcolor{gray!70}{#1. #3}%
  \fi\fi\par}
% a row of numbered boxes with values: \valrow{list}{style}
\newcommand{\valrow}[2]{%
  \foreach \v [count=\k from 1] in {#1} {
    \node[#2, minimum width=0.56cm, minimum height=0.5cm, inner sep=1pt, font=\ttfamily\scriptsize] (v\k) at ({0.6*\k},0) {\v};
  }}

\title{Cryptography}
\subtitle{Diffie-Hellman and ElGamal}
\author{Ksawery Możdżyński}
\date{}

\begin{document}

% ======================================================================
% Opening
% ======================================================================

\begin{frame}
  \titlepage
  \say{Dziś: jak dwie osoby ustalają wspólny sekret, choć ktoś słyszy każde ich słowo. I jak z tego samego pomysłu zrobić szyfr.}
\end{frame}

\begin{frame}{Agenda For Today}
  \vfill
  \begin{center}
    \begin{tikzpicture}
      % what we will do today: one icon and one sentence per point (the section names are in the bar at the top)
      \foreach \ic/\t [count=\k from 0] in {
          {\faHandshake}/{Agree on a secret key in public, without ever sending it.},
          {\faUserSecret}/{See why this is safe, and how Mallory breaks it when nobody checks who is who.},
          {\faEnvelope}/{Encrypt a message with the same idea: ElGamal.}} {
        \node[circle, fill=DeepBlue, text=white, minimum size=1.3cm, font=\LARGE] at (0,{-1.9*\k}) {\ic};
        \node[anchor=west, text width=12cm, font=\large, align=left, execute at begin node={\hyphenpenalty=10000 }] at (1.0,{-1.9*\k}) {\t};
      }
    \end{tikzpicture}
  \end{center}
  \vfill
  \say{Pierwsza połowa to wymiana klucza Diffiego-Hellmana, druga to szyfr zbudowany na tym samym pomyśle. Nazwy sekcji są na pasku u góry.}
\end{frame}

% ======================================================================
\section[A secret in public]{A secret in public}
% ======================================================================

\begin{frame}{Recall: where the shared key comes from}
  \begin{center}
    \begin{tikzpicture}[>=Stealth]
      \alice{0}{0}
      \bob{10}{0}
      \node[hwmain, minimum width=3.4cm, minimum height=1.1cm, font=\small\bfseries] (ciph) at (5,0)
        {\faBolt\ ChaCha20 / AES};
      \draw[data] (Alice) -- (ciph);
      \draw[data] (ciph) -- (Bob);
      \node[capt, below=0.15cm of ciph] {gigabytes per second, but only with the \textbf{same secret key} on both sides};
      % the missing key
      \node[draw=sAmber, very thick, dashed, rounded corners=3pt, minimum width=1.6cm, minimum height=1.0cm,
        text=sAmberText, font=\Large] (key) at (5,1.9) {\faKey\,?};
      \draw[waits] (key.west) -| (Alice.north);
      \draw[waits] (key.east) -| (Bob.north);
    \end{tikzpicture}
  \end{center}
  \vfill
  \begin{columns}[T,onlytextwidth]
    \column{0.48\textwidth}
      \stepbox{\faEnvelope\ Last lecture: RSA carries the key}{Alice picks a random key and encrypts it with Bob's \textbf{public key}.
        Only Bob's \textbf{private key} opens it. The key travels, locked.}
    \column{0.48\textwidth}
      \stepbox{\faHandshake\ Today: nobody sends it}{Alice and Bob both \textbf{compute} the same key.\\
        It never travels at all.}
  \end{columns}
  \vfill
  \say{Szyfr symetryczny: ten sam klucz szyfruje i odszyfrowuje. W poprzednim wykładzie klucz przenosiło RSA (szyfrowanie hybrydowe: kryptografia asymetryczna przenosi krótki klucz, szyfr symetryczny dane). Dziś inny sposób: klucz w ogóle nie podróżuje.}
\end{frame}

\samedot
\begin{frame}{Everybody can hear you}
  \vfill
  \begin{center}
    \begin{tikzpicture}[>=Stealth]
      \alice{0}{0}
      \bob{11}{0}
      \draw[net] (Alice.east) -- (Bob.west);
      \node[capt, fill=white] at (5.5,0.35) {\faGlobe\ internet};
      \eve{5.5}{-2.1}
      \draw[data, dashed] (5.5,-0.15) -- (Eve.north);
      \node[font=\small, text=sPinkText, anchor=west] at (6.3,-2.0) {\faAssistiveListeningSystems\ Eve hears every message};
    \end{tikzpicture}
  \end{center}
  \vfill
  \begin{center}\small They share no secret yet. Everything they say, Eve hears too.\end{center}
  \vfill
  \keymsg{Can two people agree on a secret that Eve hears in full?}
  \say{Pytanie do sali. Większość powie, że to niemożliwe. Za chwilę się okaże, że jednak można.}
\end{frame}

% ---------- the paint story ----------
% a cup of paint: \paintcup{x}{y}{colour}{label}
\newcommand{\paintcup}[4]{%
  \begin{scope}[shift={(#1,#2)}]
    \fill[#3] (-0.36,-0.42) -- (0.36,-0.42) -- (0.43,0.2) -- (-0.43,0.2) -- cycle;
    \draw[nInk, thick, line join=round] (-0.5,0.42) -- (-0.36,-0.42) -- (0.36,-0.42) -- (0.5,0.42);
    \node[font=\footnotesize, text=nInk, below] at (0,-0.44) {#4};
  \end{scope}}
\newcommand{\plus}[2]{\node[font=\large\bfseries, text=nInk] at (#1,#2) {+};}
\newcommand{\eq}[2]{\node[font=\large\bfseries, text=nInk] at (#1,#2) {=};}
% the whole picture after step #1 (1 to 4)
\newcommand{\paintscene}[1]{%
  \begin{tikzpicture}[>=Stealth, scale=0.78, transform shape]
    \useasboundingbox (-1.2,-2.6) rectangle (14.6,2.3);
    \alice{0}{-0.6}
    \bob{13.4}{-0.6}
    \eve{6.7}{1.55}
    \draw[nLight, very thick, dashed] (6.7,0.85) -- (6.7,-2.4);
    \node[capt, text=gray] at (6.7,-2.55) {public channel};
    % step 1: the public yellow
    \paintcup{1.5}{0.5}{paintY}{public}
    \paintcup{11.9}{0.5}{paintY}{public}
    \ifnum#1>1
      % step 2: each adds a secret colour
      \plus{2.35}{0.5}
      \paintcup{3.2}{0.5}{paintR}{\textcolor{sAmberText}{\faLock} secret}
      \eq{4.05}{0.5}
      \paintcup{4.9}{0.5}{paintO}{mixture}
      \plus{11.05}{0.5}
      \paintcup{10.2}{0.5}{paintB}{\textcolor{sAmberText}{\faLock} secret}
      \eq{9.35}{0.5}
      \paintcup{8.5}{0.5}{paintG}{mixture}
    \fi
    \ifnum#1>2
      % step 3: the mixtures cross the channel, Eve sees them
      \draw[flow] (5.4,0.1) to[out=-20, in=160] (11.4,-1.45);
      \draw[flow] (8.0,0.1) to[out=200, in=20] (2.0,-1.45);
      \paintcup{1.5}{-1.7}{paintG}{from Bob}
      \paintcup{11.9}{-1.7}{paintO}{from Alice}
      \node[draw=sPink, thick, fill=white, rounded corners=3pt, font=\scriptsize, text=sPinkText, anchor=west, inner sep=3pt]
        at (7.45,1.75) {Eve sees: \tikz[baseline=-0.5ex]{\fill[paintY] (0,0) circle (0.12);}
          \tikz[baseline=-0.5ex]{\fill[paintO] (0,0) circle (0.12);}
          \tikz[baseline=-0.5ex]{\fill[paintG] (0,0) circle (0.12);}};
    \fi
    \ifnum#1>3
      % step 4: each adds their own secret again
      \plus{2.35}{-1.7}
      \paintcup{3.2}{-1.7}{paintR}{\textcolor{sAmberText}{\faLock} secret}
      \eq{4.05}{-1.7}
      \paintcup{4.9}{-1.7}{paintM}{\textbf{the same}}
      \plus{11.05}{-1.7}
      \paintcup{10.2}{-1.7}{paintB}{\textcolor{sAmberText}{\faLock} secret}
      \eq{9.35}{-1.7}
      \paintcup{8.5}{-1.7}{paintM}{\textbf{the same}}
    \fi
  \end{tikzpicture}}
\newcommand{\paintsteps}[1]{%
  \footnotesize
  \stepitem{1}{#1}{Both start with a public colour: yellow.}
  \stepitem{2}{#1}{Each adds a secret colour: Alice red, Bob blue.}
  \stepitem{3}{#1}{They swap the mixtures in public.}
  \stepitem{4}{#1}{Each adds their own secret colour again.}}

\begin{frame}{Mixing paint, step 1}
  \begin{center}\paintscene{1}\end{center}
  \vfill
  \begin{center}\begin{minipage}{0.62\textwidth}\paintsteps{1}\end{minipage}\end{center}
  \vfill
  \say{Żółty jest jawny, każdy go zna, także Eve.}
\end{frame}

\samedot
\begin{frame}{Mixing paint, step 2}
  \begin{center}\paintscene{2}\end{center}
  \vfill
  \begin{center}\begin{minipage}{0.62\textwidth}\paintsteps{2}\end{minipage}\end{center}
  \vfill
  \say{Czerwony i niebieski są tajne i nigdy nie opuszczają domu.}
\end{frame}

\samedot
\begin{frame}{Mixing paint, step 3}
  \begin{center}\paintscene{3}\end{center}
  \vfill
  \begin{center}\begin{minipage}{0.62\textwidth}\paintsteps{3}\end{minipage}\end{center}
  \vfill
  \say{Eve widzi obie mieszanki: żółto-czerwoną i żółto-niebieską.}
\end{frame}

\samedot
\begin{frame}{Mixing paint, step 4}
  \begin{center}\paintscene{4}\end{center}
  \vspace{-0.15cm}
  \begin{center}\begin{minipage}{0.62\textwidth}\paintsteps{4}
    \vspace{2pt}\phantom{$\blacktriangleright$}\ \ Both now have yellow + red + blue: \textbf{the same colour}.\end{minipage}\end{center}
  \vfill
  \keymsg{Mixing is easy, unmixing is hard.}
  \say{Kolejność mieszania nie ma znaczenia, więc obie strony mają ten sam kolor. Eve ma tylko mieszanki, a rozdzielić farby się nie da.}
\end{frame}

\begin{frame}{Diffie and Hellman, 1976}
  \vfill
  \begin{columns}[c,onlytextwidth]
    \column{0.42\textwidth}
      \begin{center}
      \begin{tikzpicture}
        % the first page of the paper
        \draw[nLine, thick, fill=white, drop shadow={opacity=0.25}] (0,0) rectangle (4.2,5.2);
        \node[font=\small\bfseries, text=nInk, text width=3.6cm, align=center] at (2.1,4.4) {New Directions in Cryptography};
        \node[font=\scriptsize, text=nInk] at (2.1,3.6) {Whitfield Diffie};
        \node[font=\scriptsize, text=nInk] at (2.1,3.25) {Martin E. Hellman};
        \foreach \y in {2.7,2.45,...,0.4} \draw[nLight, thick] (0.4,\y) -- (3.8,\y);
        \node[fill=DeepBlue, text=white, font=\small\bfseries, rounded corners=2pt] at (3.4,0.3) {1976};
      \end{tikzpicture}
      \end{center}
    \column{0.55\textwidth}
      Whitfield Diffie and Martin Hellman,\\ \textbf{``New Directions in Cryptography''}, 1976.
      \vspace{0.4cm}

      The first public way to \textbf{agree on a secret over an open channel}.
      \vspace{0.4cm}

      \textcolor{AccentOrange}{\Large\faTrophy}\ \ Turing Award, 2015.
  \end{columns}
  \vfill
  % CHECKED 2026-11-25, read: Hellman, "An overview of public key cryptography", IEEE Communications Magazine, 50th Anniversary
  %   issue, May 2002 (www.cs.miami.edu/home/burt/manuscripts/crypto_for_intelligence/hellman.pdf): "hence should be called
  %   Diffie-Hellman-Merkle key exchange". The other years below: secondary sources only.
  % CHECKED (secondary sources: Wikipedia "Merkle's Puzzles", "Malcolm J. Williamson"; schneier.com 1998): the years below.
  %   Turing Award 2015 (announced March 2016), amturing.acm.org/award_winners/diffie_8371646.cfm;
  %   Merkle's puzzles: a CS244 project at Berkeley in 1974, published in CACM 21(4), 1978; Hellman proposed the name
  %   "Diffie-Hellman-Merkle" in IEEE Communications Magazine, May 2002 ("An overview of public key cryptography");
  %   GCHQ: Malcolm Williamson, 1974 to 1976, made public by GCHQ in 1997
  \say{Ralph Merkle miał podobny pomysł już w 1974 roku (zagadki Merkle'a, projekt na zajęciach w Berkeley, opublikowany dopiero w 1978); Hellman w 2002 zaproponował nazwę Diffie-Hellman-Merkle. Słowo ,,public'' jest ważne: w brytyjskim GCHQ Malcolm Williamson wymyślił tę samą wymianę w latach 1974 do 1976, ale to było tajne do 1997 roku (tak samo jak RSA w wykładzie 3). Nagroda Turinga to najważniejsza nagroda w informatyce; ta jest za rok 2015, ogłoszona w marcu 2016.}
\end{frame}

% ======================================================================
\section[Diffie-Hellman]{Diffie-Hellman}
% ======================================================================

\begin{frame}{Recall: powers modulo p}
  \vfill
  \begin{center}
    \LARGE $5^2 \bmod 23 \;=\; 25 \bmod 23 \;=\; \textcolor{AccentOrange}{\mathbf{2}}$
  \end{center}
  \vfill
  \begin{center}
    \normalsize In Python:\quad \pyprompt{pow(5, 2, 23)}{2}
  \end{center}
  \vfill
  \begin{center}\small Fast even when every number has \textbf{600 digits}.\end{center}
  \vfill
  \curious{diffie-hellman/\#5}
  % CHECKED (python3): pow(5, 2, 23) = 2
  \say{Mod to reszta z dzielenia. pow z trzema argumentami liczy potęgę modulo bez liczenia ogromnej potęgi. Jak to robi (square-and-multiply), pokazuje strona z linku, krok po kroku po bitach wykładnika, na liczbach z tego wykładu.}
\end{frame}

\begin{frame}{The public numbers: p and g}
  \vfill
  \begin{center}
    \begin{tikzpicture}
      \node[pub, minimum width=4.2cm, minimum height=2.2cm] (p) at (0,0) {};
      \node[font=\Huge\bfseries, text=nInk] at (0,0.3) {$p = 23$};
      \node[capt] at (0,-0.65) {a prime};
      \node[pub, minimum width=4.2cm, minimum height=2.2cm] (g) at (5.5,0) {};
      \node[font=\Huge\bfseries, text=nInk] at (5.5,0.3) {$g = 5$};
      \node[capt] at (5.5,-0.65) {a number smaller than $p$};
      \foreach \n in {p,g} \node[fill=nInk, text=white, font=\tiny\bfseries, rounded corners=2pt, anchor=south east]
        at (\n.north east) {\faEye\ public};
    \end{tikzpicture}
  \end{center}
  \vfill
  \begin{center}
    Both are \textbf{public}: everybody may know them, Eve too.\\[2pt]
    They play the role of the \tikz[baseline=-0.6ex]{\fill[paintY] (0,0) circle (0.13);} \textbf{yellow paint}.
  \end{center}
  \vfill
  \say{Liczba pierwsza dzieli się tylko przez 1 i siebie. Dlaczego pierwsza: wtedy każda liczba od 1 do p - 1 ma odwrotność modulo p, a potęgi g mogą obejść je wszystkie; przy p złożonym niektóre liczby nie mają odwrotności i cykle rozpadają się na krótsze kawałki. W prawdziwej wymianie p ma 2048 bitów; teraz małe liczby, żeby liczyć na kartce.}
\end{frame}

% CHECKED (python3): [pow(5, e, 23) for e in range(1, 23)]
\samedot
\begin{frame}{Why this g?}
  \small Powers of 5 mod 23, for exponents 1 to 22:
  \begin{center}
    \begin{tikzpicture}
      \foreach \e in {1,...,22} \node[font=\tiny, text=gray] at ({0.6*\e},0.45) {\e};
      \valrow{5,2,10,4,20,8,17,16,11,9,22,18,21,13,19,3,15,6,7,12,14,1}{pub}
      \node[rowlabel] at (0.2,0.45) {exponent};
      \node[rowlabel] at (0.2,0) {$5^e \bmod 23$};
    \end{tikzpicture}
  \end{center}
  \vspace{-0.3cm}
  \begin{center}
    \begin{tikzpicture}
      \begin{axis}[width=12cm, height=3.9cm, xmin=0, xmax=23, ymin=0, ymax=23.5,
          xtick={1,5,10,15,20,22}, ytick={1,5,10,15,20,22}, tick label style={font=\tiny},
          axis lines=left, axis line style={nLine}, xlabel={\scriptsize exponent $e$},
          ylabel={\scriptsize $5^e \bmod 23$}, ymajorgrids, grid style={nLight}]
        \addplot[only marks, mark=*, mark size=2pt, color=DeepBlue] coordinates {
          (1,5)(2,2)(3,10)(4,4)(5,20)(6,8)(7,17)(8,16)(9,11)(10,9)(11,22)
          (12,18)(13,21)(14,13)(15,19)(16,3)(17,15)(18,6)(19,7)(20,12)(21,14)(22,1)};
      \end{axis}
    \end{tikzpicture}
  \end{center}
  \vspace{-0.25cm}
  % CHECKED (python3): pow(5, 23, 23) = 5
  \begin{center}\footnotesize Every number from 1 to 22 appears \textbf{exactly once}. Then $5^{23} \bmod 23 = 5$ again:
    the powers go round in a \textbf{cycle} of 22.\end{center}
  \keymsg{The cycle of $g$ is what Eve has to search, so it must be long.}
  \say{Taką liczbę g nazywamy generatorem. Liczy się długość cyklu: ile różnych wartości dają potęgi g, zanim zaczną się powtarzać. Tyle możliwości musiałaby sprawdzić Eve. Widać też, że wyniki skaczą bez żadnego porządku; to wróci w sekcji 3.}
\end{frame}

\samedot
\begin{frame}{A bad g}
  \vfill
  \normalsize Powers of 22 mod 23:
  \begin{center}
    \begin{tikzpicture}[scale=1.35, transform shape]
      % CHECKED (python3): [pow(22, e, 23) for e in range(1, 17)] = 22, 1, 22, 1, ...
      \foreach \e in {1,...,16} \node[font=\tiny, text=gray] at ({0.6*\e},0.45) {\e};
      \valrow{22,1,22,1,22,1,22,1,22,1,22,1,22,1,22,1}{pub}
      \foreach \k in {3,...,16} \node[cellfree, minimum width=0.56cm, minimum height=0.5cm, fill=AccentOrange!12,
        draw=AccentOrange, font=\ttfamily\scriptsize, solid] at (v\k) {};
      \foreach \k/\v in {3/22,4/1,5/22,6/1,7/22,8/1,9/22,10/1,11/22,12/1,13/22,14/1,15/22,16/1}
        \node[font=\ttfamily\scriptsize, text=nInk] at (v\k) {\v};
      \draw[-{Stealth}, very thick, AccentOrange] (v2.south) to[out=-40, in=-140] node[below, font=\scriptsize, text=AccentOrange] {back to the start}
        ([xshift=-2pt]v3.south);
    \end{tikzpicture}
  \end{center}
  \vfill
  \begin{center}\large Only \textbf{2 different values}: $A$ is 22 or 1, whatever Alice's secret $a$ is.\\[4pt]
    Eve guesses the shared secret on the first or the second try.\end{center}
  \vfill
  \curious{diffie-hellman/\#2}
  \keymsg{A short cycle is a disaster: $g$ must have a long cycle.}
  \say{22 to p - 1, czyli -1 modulo 23: jego potęgi to na zmianę 22 i 1. Taki cykl to katastrofa. Wybór g należy do twórców standardu, nie do użytkownika. Ta liczba 22 wróci w sekcji 5: wtedy wybierze ją Mallory, nie standard. Dlaczego prawdziwe standardy biorą cykl o połowę krótszy niż 22: w sekcji 3. W symulatorze z linku potęgi g chodzą po zegarze; wpisz g = 22.}
\end{frame}

% ---------- the exchange, one picture in five steps ----------
\newcommand{\dhrow}[3]{% \dhrow{step number}{current step}{y}: highlight when current
  \ifnum#1=#2\relax \fill[rowhl] (-0.3,{#3-0.36}) rectangle (14.3,{#3+0.36}); \fi}
\newcommand{\dhlabel}[4]{% \dhlabel{number}{current}{y}{text}
  \ifnum#1=#2\relax
    \node[anchor=west, font=\small\bfseries, text=AccentOrange] at (-0.2,#3) {$\blacktriangleright$ #1. #4};
  \else
    \node[anchor=west, font=\small, text=nInk] at (0.1,#3) {#1. #4};
  \fi}
\newcommand{\dhscene}[1]{%
  \begin{tikzpicture}[>=Stealth, scale=0.97, transform shape]
    \useasboundingbox (-0.4,-4.25) rectangle (14.4,1.55);
    \node[circle, fill=sIndigo, text=white, minimum size=0.75cm, inner sep=0pt] (an) at (4.6,1.0) {\faUser};
    \node[font=\small\bfseries, text=sIndigo, right=2pt of an] {Alice};
    \node[circle, fill=sTeal, text=white, minimum size=0.75cm, inner sep=0pt] (bt) at (11.6,1.0) {\faUser};
    \node[font=\small\bfseries, text=sTealText, right=2pt of bt] {Bob};
    % Eve's notebook over the channel
    \node[circle, fill=sPink, text=white, minimum size=0.75cm, inner sep=0pt] (ew) at (8.25,1.0) {\faUserSecret};
    \ifnum#1>2
      \node[draw=sPink, thick, fill=white, rounded corners=2pt, font=\scriptsize, text=sPinkText, anchor=west, align=left]
        at (ew.east) {\faBook\ 23, 5, 8, 19};
    \else
      \node[draw=sPink, thick, fill=white, rounded corners=2pt, font=\scriptsize, text=sPinkText, anchor=west, align=left]
        at (ew.east) {\faBook\ 23, 5};
    \fi
    \draw[nLight, very thick, dashed] (8.25,0.5) -- (8.25,-3.55);
    \dhrow{1}{#1}{0}  \dhrow{2}{#1}{-0.8}  \dhrow{3}{#1}{-1.6}  \dhrow{4}{#1}{-2.4}  \dhrow{5}{#1}{-3.2}
    \dhlabel{1}{#1}{0}{Secrets}
    \node[sec] at (5.1,0) {$a = 6$};
    \node[sec] at (12.1,0) {$b = 15$};
    \ifnum#1>1
      \dhlabel{2}{#1}{-0.8}{Public values}
      \node[pub] at (5.1,-0.8) {$A = 5^{6} \bmod 23 = 8$};
      \node[pub] at (12.1,-0.8) {$B = 5^{15} \bmod 23 = 19$};
    \fi
    \ifnum#1>2
      \dhlabel{3}{#1}{-1.6}{Send}
      \draw[netarrow] (5.6,-1.42) -- node[above=3pt, font=\scriptsize, text=nInk, fill=white, inner sep=1pt] {sends $A = 8$} (11.0,-1.45);
      \draw[netarrow] (11.0,-1.8) -- node[below=3pt, font=\scriptsize, text=nInk, fill=white, inner sep=1pt] {sends $B = 19$} (5.6,-1.8);
    \fi
    \ifnum#1>3
      \dhlabel{4}{#1}{-2.4}{Alice computes}
      \node[sec] at (5.1,-2.4) {$B^{a} = 19^{6} \bmod 23 = \mathbf{2}$};
    \fi
    \ifnum#1>4
      \dhlabel{5}{#1}{-3.2}{Bob computes}
      \node[sec] at (12.1,-3.2) {$A^{b} = 8^{15} \bmod 23 = \mathbf{2}$};
    \fi
    \node[pub, font=\small] at (5.1,-4.0) {$p = 23,\; g = 5$\quad (public)};
  \end{tikzpicture}}

\begin{frame}{The exchange, step 1}
  \vfill
  \begin{center}\dhscene{1}\end{center}
  \vfill
  \say{Alice i Bob wybierają tajne liczby. Eve ich nie zna.}
\end{frame}

\samedot
\begin{frame}{The exchange, step 2}
  \vfill
  \begin{center}\dhscene{2}\end{center}
  \vfill
  \say{Każdy liczy g do swojej tajnej potęgi. To są mieszanki farb.}
\end{frame}

\samedot
\begin{frame}{The exchange, step 3}
  \vfill
  \begin{center}\dhscene{3}\end{center}
  \vfill
  \say{Eve widzi 8 i 19. Zapisuje je.}
\end{frame}

\samedot
\begin{frame}{The exchange, step 4}
  \vfill
  \begin{center}\dhscene{4}\end{center}
  \vfill
  \say{Alice bierze liczbę od Boba i podnosi do swojej tajnej potęgi.}
\end{frame}

% CHECKED (python3): pow(5,6,23) = 8, pow(5,15,23) = 19, pow(19,6,23) = 2, pow(8,15,23) = 2
\samedot
\begin{frame}{The exchange, step 5}
  \vfill
  \begin{center}\dhscene{5}\end{center}
  \vfill
  \keymsg{Both have 2. Eve saw only 23, 5, 8 and 19.}
  \say{Bob robi to samo z liczbą od Alice. Ten sam wynik: 2.}
\end{frame}

\samedot
\begin{frame}{Why are they equal?}
  \vfill
  \begin{center}
    \begin{tikzpicture}
      \node[anchor=east, font=\large\bfseries, text=sIndigo] at (0,0.7) {Alice:};
      \node[anchor=west, font=\Large] at (0.2,0.7) {$B^{a} \;=\; (g^{b})^{a} \;=\; g^{\,b a}$};
      \node[anchor=east, font=\large\bfseries, text=sTealText] at (0,-0.7) {Bob:};
      \node[anchor=west, font=\Large] at (0.2,-0.7) {$A^{b} \;=\; (g^{a})^{b} \;=\; g^{\,a b}$};
      \node[sec, font=\Large, minimum height=2.5cm, minimum width=3.2cm] (r) at (8.8,0) {$g^{ab} \bmod p$};
      \draw[data] (5.6,0.7) -- (r.west |- 0,0.7);
      \draw[data] (5.6,-0.7) -- (r.west |- 0,-0.7);
    \end{tikzpicture}
  \end{center}
  \vfill
  % CHECKED (python3): 6 * 15 = 90, pow(5, 90, 23) = 2
  \begin{center}$a\,b = b\,a$. In our exchange: $6 \times 15 = 90$, and $5^{90} \bmod 23 = 2$.\end{center}
  \vfill
  \curious{diffie-hellman/\#1}
  \keymsg{Both compute $g^{(ab)} \bmod p$, and nobody ever sent it.}
  \say{Zwykłe prawo potęg, działa też modulo p, bo (x mod p)\textasciicircum a mod p = x\textasciicircum a mod p, jak w RSA: można redukować modulo p w dowolnym momencie. To jest zamiana kolejności mieszania farb. W symulatorze z linku ta sama wymiana jako czarna skrzynka: można zmienić a, b, p i g i patrzeć, jak oba sekrety zmieniają się razem.}
\end{frame}

% the table for task 1: #1 = 0 empty, 1 filled
\newcommand{\taskonetable}[1]{%
  \renewcommand{\arraystretch}{1.5}
  \begin{tabular}{lll}
    \toprule
    & \textbf{computation} & \textbf{result} \\
    \midrule
    $A$ & \ifnum#1=1 $5^3 \bmod 23 = 125 \bmod 23$\else \rule{4.2cm}{0pt}\fi & \ifnum#1=1 $10$\fi \\
    $B$ & \ifnum#1=1 $5^5 \bmod 23 = 3125 \bmod 23$\fi & \ifnum#1=1 $20$\fi \\
    \textcolor{sIndigo}{\textbf{Alice}}: $B^a$ & \ifnum#1=1 $20^3 \bmod 23 = 8000 \bmod 23$\fi &
      \ifnum#1=1 \colorbox{AccentOrange!20}{\textbf{19}}\fi \\
    \textcolor{sTealText}{\textbf{Bob}}: $A^b$ & \ifnum#1=1 $10^5 \bmod 23 = 100000 \bmod 23$\fi &
      \ifnum#1=1 \colorbox{AccentOrange!20}{\textbf{19}}\fi \\
    \bottomrule
  \end{tabular}}

\begin{frame}{Task 1: an exchange on paper}
  \begin{center}\large $p = 23$,\quad $g = 5$,\quad Alice's secret $a = 3$,\quad Bob's secret $b = 5$\end{center}
  \begin{columns}[T,onlytextwidth]
    \column{0.42\textwidth}
      \small
      \begin{enumerate}\setlength{\itemsep}{6pt}
        \item Compute $A$ and $B$.
        \item Compute the shared secret as Alice, and again as Bob.
        \item Check that both are the same.
      \end{enumerate}
    \column{0.56\textwidth}
      \footnotesize\taskonetable{0}
  \end{columns}
  \vfill
  \onpaper{about 5 minutes}
  \vfill
  \say{Około 5 minut, na kartce. Potęgi są małe: 5\textasciicircum 3 = 125, 5\textasciicircum 5 = 3125. Wolniejszym podpowiedzieć: najpierw A i B, potem każdy podnosi liczbę od drugiej strony do swojej tajnej potęgi.}
\end{frame}

% CHECKED (python3): pow(5,3,23) = 10, pow(5,5,23) = 20, pow(20,3,23) = 19, pow(10,5,23) = 19; 20**3 = 8000, 10**5 = 100000
\samedot
\begin{frame}{Task 1: what we saw}
  \vfill
  \begin{center}\small
    \renewcommand{\arraystretch}{1.6}
    \begin{tabular}{rl}
      $A$ =& $5^3 \bmod 23 = 125 \bmod 23 = 10$ \\
      $B$ =& $5^5 \bmod 23 = 3125 \bmod 23 = 20$ \\
      \textcolor{sIndigo}{\textbf{Alice}}:\ \ $B^a$ =& $20^3 \bmod 23 = 8000 \bmod 23 = \colorbox{AccentOrange!20}{\textbf{19}}$ \\
      \textcolor{sTealText}{\textbf{Bob}}:\ \ $A^b$ =& $10^5 \bmod 23 = 100000 \bmod 23 = \colorbox{AccentOrange!20}{\textbf{19}}$ \\
    \end{tabular}
  \end{center}
  \vfill
  \keymsg{Both get 19. Eve saw only 23, 5, 10 and 20.}
  \say{Sprawdzone w Pythonie. 8000 = 347 · 23 + 19, 100000 = 4347 · 23 + 19.}
\end{frame}

% CHECKED (python3): the prime of group 14 from RFC 3526 has 2048 bits and 617 decimal digits
% CHECKED: secrets of 256 bits: RFC 3526, section 8, gives 220 to 320 bits of exponent for group 14;
%   NIST SP 800-56A Rev. 3 asks for at least 2 * 112 = 224 bits at 112 bits of security; 256 bits give 2^128 steps for the
%   square-root methods (slide "From 28 bits to 256 bits"). TLS 1.3 uses the groups ffdhe2048 ... of RFC 7919 (section 4.2.7 of RFC 8446)
\begin{frame}{Real sizes}
  \vfill
  \begin{center}
    \renewcommand{\arraystretch}{1.45}
    \begin{tabular}{@{}lll@{}}
      \tikz[baseline=(n.base)]\node[pub] (n) {$p$}; & \textbf{2048 bits} (617 digits) & a fixed public prime from the standard \textbf{RFC 3526} \\
      \tikz[baseline=(n.base)]\node[pub] (n) {$g$}; & 2 & from the same standard \\
      \tikz[baseline=(n.base)]\node[sec] (n) {$a, b$}; & at least \textbf{256 bits} & random and \textbf{secret} \ \textcolor{gray}{(why 256: section 3)} \\
    \end{tabular}
  \end{center}
  \vspace{0.1cm}
  \begin{center}\small A fixed public pair $p$, $g$ that everybody may use is called a \textbf{group}.\end{center}
  \begin{beamercolorbox}[rounded=true,sep=0.15cm,wd=\linewidth]{block body}
    \scriptsize\ttfamily\color{nInk} $p$ = \seqsplit{32317006071311007300338913926423828248817941241140239112842009751400741706634354222619689417363569347117901737909704191754605873209195028853758986185622153212175412514901774520270235796078236248884246189477587641105928646099411723245426622522193230540919037680524235519125679715870117001058055877651038861847280257976054903569732561526167081339361799541336476559160368317896729073178384589680639671900977202194168647225871031411336429319536193471636533209717077448227988588565369208645296636077250268955505928362751121174096972998068410554359584866583291642136218231078990999448652468262416972035911852507045361090559}
  \end{beamercolorbox}
  \vfill
  \say{RFC to oficjalny dokument standardu internetowego; RFC 3526 podaje kilka gotowych par p i g (grupa 14 ma 2048 bitów). Słowo grupa ma w matematyce szersze znaczenie; tu wystarczy: gotowa para p i g, wspólna dla wszystkich. Gdyby ktoś pytał: TLS 1.3 nie bierze grup z RFC 3526, tylko nazwane grupy ffdhe2048, ffdhe3072 i dalsze z RFC 7919; inne liczby pierwsze, ta sama budowa (bezpieczna liczba pierwsza, g = 2).}
\end{frame}

\begin{frame}{Task 2: a real exchange}
  \vfill
  \taskgoal{one Diffie-Hellman exchange with the real 2048-bit $p$, and how long it takes.}
  \taskready{$p$ and $g = 2$ from RFC 3526, and the timer, in \texttt{dh.py}.}
  \vfill
  \opentask{Task2}{about 6 minutes}
  \vfill
  \say{Około 6 minut. Liczby losowe: secrets.randbits(256). Moduł secrets daje liczby losowe dobre do kryptografii, w przeciwieństwie do random. Każde zadanie jest samodzielne, nic nie trzeba instalować. Instrukcja i TODO są w README na GitHubie.}
\end{frame}

% CHECKED: python3 Task2/solution/dh.py on the author's laptop (Intel Core i5-8265U), 2026-11-17, five runs:
%          11.3, 11.6, 11.3, 11.4, 11.7 ms; the run shown below is the fourth one (the same as in Task2/README.md);
%          the 20 digits of the secret are new on every run
\samedot
\begin{frame}{Task 2: what we saw}
  \vfill
  \begin{center}
    \begin{minipage}{0.6\textwidth}
      \stepbox{The output}{\ttfamily equal: True\\ secret: 11974210656905040766 ...\\ took 11.4 ms}
    \end{minipage}\\[0.4cm]
    Both sides get the same secret of up to 617 digits, a new one on every run.\\[2pt]
    The whole exchange takes about \textbf{11 ms} in plain Python (on this laptop, Intel Core i5-8265U).
  \end{center}
  \vfill
  \keymsg{A few lines of Python, and two strangers share a 2048-bit secret.}
  \say{Sprawdzone na tym laptopie (i5-8265U), pięć uruchomień: od 11,3 do 11,7 ms. Rozwiązanie: a i b z secrets.randbits(256), A = pow(G, a, P), B = pow(G, b, P), s\_alice = pow(B, a, P), s\_bob = pow(A, b, P). Cyfry sekretu są inne przy każdym uruchomieniu. OpenSSL robi całą stronę wymiany w około 0,3 ms.}
\end{frame}

% ======================================================================
\section[Why is it safe?]{Why is it safe?}
% ======================================================================

\begin{frame}{What Eve knows}
  \vfill
  \begin{columns}[c,onlytextwidth]
    \column{0.4\textwidth}
      \begin{center}
      \begin{tikzpicture}
        % Eve's notebook
        \draw[sPink, thick, fill=white, rounded corners=3pt] (0,0) rectangle (4.4,4.3);
        \fill[sPink] (0,3.75) rectangle (4.4,4.3);
        \node[text=white, font=\small\bfseries] at (2.2,4.02) {\faUserSecret\ Eve's notebook};
        \foreach \y in {0.4,0.95,...,3.3} \draw[nLight] (0.25,\y) -- (4.15,\y);
        \node[anchor=west, font=\small, text=nInk] at (0.3,3.3) {$p = 23,\; g = 5$};
        \node[anchor=west, font=\small, text=nInk] at (0.3,2.6) {$A = 8$};
        \node[anchor=west, font=\small, text=nInk] at (0.3,1.9) {$B = 19$};
        \node[anchor=west, font=\small, text=gray] at (0.3,1.2) {$a = \,?$};
        \node[anchor=west, font=\small, text=gray] at (2.3,1.2) {$b = \,?$};
        \draw[aWrite, very thick] (0.3,1.0) -- (1.4,1.4) (2.3,1.0) -- (3.4,1.4);
      \end{tikzpicture}
      \end{center}
    \column{0.56\textwidth}
      \large Eve has everything that was sent.
      \vspace{0.4cm}

      To compute the secret $g^{ab}$, she needs \textcolor{sAmberText}{\textbf{$a$ or $b$}}.
      \vspace{0.4cm}

      % CHECKED (python3): 8 * 19 % 23 = 14, pow(5, 21, 23) = 14, pow(5, 90, 23) = 2
      \small Eve can multiply: $A \times B = g^{a+b}$, not $g^{ab}$:\\ $8 \times 19 \bmod 23 = 14$, not 2.
      \vspace{0.4cm}

      \normalsize Can she get $a$ back from $A = 8$?
  \end{columns}
  \vfill
  \say{Eve ma wszystko poza tajnymi liczbami. Pomnożenie A i B dodaje wykładniki, a potrzebny jest ich iloczyn: 5\textasciicircum 21 mod 23 = 14, a nie 2. W języku farb: zmieszanie obu mieszanek daje kolor z dwoma porcjami żółtego, a nie ten wspólny. Czy da się odzyskać a z A?}
\end{frame}

\samedot
\begin{frame}{The discrete logarithm}
  \bigdef{\textbf{Discrete logarithm:} given $p$, $g$ and $A = g^{a} \bmod p$, find $a$.}
  \begin{columns}[T,onlytextwidth]
    \column{0.49\textwidth}
      \centering
      \begin{tikzpicture}
        \begin{axis}[width=6.6cm, height=3.5cm, xmin=0, xmax=6.5, ymin=0, ymax=17000, axis lines=left,
            axis line style={nLine}, tick label style={font=\tiny}, xtick={0,...,6},
            ytick={0,5000,10000,15625}, yticklabel style={/pgf/number format/fixed},
            scaled y ticks=false, title={\scriptsize over normal numbers: $5^a$}]
          \addplot[DeepBlue, very thick, domain=0:6.2, samples=60] {5^x};
          \addplot[only marks, mark=*, mark size=1.8pt, DeepBlue] coordinates {(1,5)(2,25)(3,125)(4,625)(5,3125)(6,15625)};
          \addplot[only marks, mark=*, mark size=3pt, AccentOrange] coordinates {(6,15625)};
        \end{axis}
      \end{tikzpicture}\\
      \footnotesize $5^a = 15625 \;\rightarrow\; a = 6$, \textbf{easy}: the powers grow steadily
    \column{0.49\textwidth}
      \centering
      \begin{tikzpicture}
        \begin{axis}[width=6.6cm, height=3.5cm, xmin=0, xmax=23, ymin=0, ymax=23.5, axis lines=left,
            axis line style={nLine}, tick label style={font=\tiny}, xtick={1,6,11,16,22}, ytick={1,8,15,22},
            title={\scriptsize modulo 23: $5^a \bmod 23$}]
          \addplot[only marks, mark=*, mark size=1.8pt, DeepBlue] coordinates {
            (1,5)(2,2)(3,10)(4,4)(5,20)(6,8)(7,17)(8,16)(9,11)(10,9)(11,22)
            (12,18)(13,21)(14,13)(15,19)(16,3)(17,15)(18,6)(19,7)(20,12)(21,14)(22,1)};
          \addplot[only marks, mark=*, mark size=3pt, AccentOrange] coordinates {(6,8)};
        \end{axis}
      \end{tikzpicture}\\
      \footnotesize $5^a \bmod 23 = 8 \;\rightarrow$ the powers \textbf{jump around}:\\ 5, 2, 10, 4, 20, 8, \dots
  \end{columns}
  \vfill
  \curious{diffie-hellman/\#2}
  \keymsg{Powers modulo p are easy to compute and hard to undo.}
  \say{Zwykły logarytm da się policzyć, bo potęgi rosną po kolei. Modulo p nie ma porządku, więc trzeba zgadywać. W symulatorze z linku potęgi g na tarczy zegara i Eve, która zgaduje wykładnik po kolei.}
\end{frame}

% CHECKED (python3): pow(8, 11, 23) = 1, pow(19, 11, 23) = 22; pow(5, e, 23) ** 11 % 23 == 1 exactly for the 11 even e in 1..22
\begin{frame}{A cycle of 22 gives away one bit}
  \small Powers of 5 mod 23 again. The cycle has $22 = 2 \times 11$ values. Raise any of them to the power 11:
  \begin{center}
    \begin{tikzpicture}
      \foreach \e in {1,...,22} \node[font=\tiny, text=gray] at ({0.6*\e},0.45) {\e};
      \valrow{5,2,10,4,20,8,17,16,11,9,22,18,21,13,19,3,15,6,7,12,14,1}{pub}
      \foreach \k in {2,4,...,22} \node[draw=aExec, very thick, rounded corners=2pt, minimum width=0.56cm, minimum height=0.5cm] at (v\k) {};
      \node[rowlabel] at (0.2,0.45) {exponent $e$};
      \node[rowlabel] at (0.2,0) {$5^e \bmod 23$};
      \node[font=\scriptsize, text=aExec, anchor=west] at (0.3,-0.55) {\faSquare[regular]\ $e$ even: the value to the power 11 gives \textbf{1}};
      \node[font=\scriptsize, text=nInk, anchor=west] at (7.3,-0.55) {$e$ odd: it gives \textbf{22}};
    \end{tikzpicture}
  \end{center}
  \begin{columns}[T,onlytextwidth]
    \column{0.48\textwidth}
      \stepbox{\faUserSecret\ Eve, with Alice's $A = 8$}{$8^{11} \bmod 23 = \mathbf{1}$, so $a$ is \textbf{even}.\\ \textcolor{gray}{(Indeed $a = 6$.)}}
    \column{0.48\textwidth}
      \stepbox{\faUserSecret\ Eve, with Bob's $B = 19$}{$19^{11} \bmod 23 = \mathbf{22}$, so $b$ is \textbf{odd}.\\ \textcolor{gray}{(Indeed $b = 15$.)}}
  \end{columns}
  \vfill
  \begin{center}\small Why: if $a = 2t$, then $A^{11} = 5^{22t} = 1$, because 22 steps go once round the cycle.\end{center}
  \keymsg{Every small factor of the cycle length gives Eve a piece of $a$ for free.}
  \say{Eve nie liczy logarytmu, tylko jedną potęgę: A do potęgi 11. Wynik 1 znaczy, że a jest parzyste. To jeden bit za darmo. Ogólnie (metoda Pohliga-Hellmana): jeśli długość cyklu rozkłada się na małe czynniki, Eve liczy osobno kawałek a dla każdego czynnika i skleja je; logarytm jest tak trudny jak dla największego czynnika pierwszego długości cyklu. Przy 22 = 2 · 11 zostaje do szukania tylko cykl 11. Pytanie do sali: czy Eve może tak zgadnąć parzystość, gdy cykl ma 11 wartości? Następny slajd.}
\end{frame}

% CHECKED (python3): [pow(2, e, 23) for e in range(1, 12)] = 2, 4, 8, 16, 9, 18, 13, 3, 6, 12, 1; 23 = 2 * 11 + 1;
%   RFC 3526 group 14: q = (P - 1)//2 has 2047 bits and passes 40 rounds of Miller-Rabin, pow(2, q, P) == 1
\samedot
\begin{frame}{Why a cycle of prime length?}
  \small Powers of 2 mod 23: a cycle of \textbf{11} values, and 11 is a prime.
  \begin{center}
    \begin{tikzpicture}[scale=1.3, transform shape]
      \foreach \e in {1,...,11} \node[font=\tiny, text=gray] at ({0.6*\e},0.45) {\e};
      \valrow{2,4,8,16,9,18,13,3,6,12,1}{pub}
      \node[cellfree, minimum width=0.56cm, minimum height=0.5cm] at (7.2,0) {2};
      \node[cellfree, minimum width=0.56cm, minimum height=0.5cm] at (7.8,0) {4};
      \node[font=\scriptsize, text=gray] at (8.4,0) {\dots};
      \draw[-{Stealth}, very thick, AccentOrange] (v11.south) to[out=-40, in=-140] node[below, font=\tiny, text=AccentOrange] {again} (7.15,-0.27);
    \end{tikzpicture}
  \end{center}
  \vspace{-0.1cm}
  \begin{itemize}\setlength{\itemsep}{4pt}
    \item 11 has no smaller factors, so there are \textbf{no small pieces} to split off: Eve has to search the whole cycle.
    \item It costs only 1 bit: 11 values instead of 22.
    \item The standard does exactly this. $p = 2q + 1$ with $q$ also prime (a \textbf{safe prime}), and $g$ with a cycle of exactly $q$.
      Here $23 = 2 \times 11 + 1$. In RFC 3526, $g = 2$ and $q$ is a prime of 2047 bits.
  \end{itemize}
  \vfill
  \curious{diffie-hellman/\#2}
  \keymsg{Real groups use a cycle of prime length $q$: long, and with no small pieces.}
  \say{Dlatego pół cyklu to nie wada, tylko cel: pełny cykl ma długość p - 1, a to liczba parzysta, więc zawsze ma czynnik 2. Bezpieczna liczba pierwsza p = 2q + 1 daje cykl długości q, liczby pierwszej. Ten sam pomysł da nam w sekcji 5 sprawdzenie wartości od drugiej strony. W zadaniu 3 wszystkie cztery p są takimi bezpiecznymi liczbami pierwszymi, a g = 4 ma cykl q. Dla ciekawych: w pełnej grupie (g = 5) A zdradza też, czy wynik jest ,,kwadratem'' modulo p, co psuje szyfrowanie ElGamal z sekcji 6; w grupie rzędu pierwszego tego problemu nie ma. W symulatorze z linku przycisk g = 2 pokazuje ten cykl na tarczy.}
\end{frame}

\begin{frame}{Task 3: find a}
  \vfill
  \taskgoal{find $a$ from $A$ by trying $a = 1, 2, 3, \dots$, then estimate how long a 256-bit $a$ would take.}
  \taskready{four problems ($p$ of 16, 20, 24 and 28 bits) and the timer, in \texttt{dlog.py}.}
  \begin{center}
    \small\textcolor{AccentOrange}{\faLightbulb}\ \textbf{Predict first:} each $p$ is 4 bits longer. How many times longer will it take?
  \end{center}
  \vfill
  \opentask{Task3}{about 8 minutes}
  \vfill
  \say{Około 8 minut. Zapisać przewidywania na tablicy. Ostatni przypadek liczy się kilka sekund. Nie liczymy pow od nowa dla każdego a: wystarczy pomnożyć poprzednią wartość przez g. Drugi krok: własna prędkość (a / seconds prób na sekundę) i ile lat zajęłoby 2\textasciicircum 256 prób.}
\end{frame}

% CHECKED: python3 Task3/solution/dlog.py on the author's laptop (Intel Core i5-8265U), 2026-11-17, three runs;
%   the second run is shown (the same as in Task3/README.md): 0.002, 0.035, 0.336, 5.836 s,
%   "10374805 tries per second, a of 256 bits: about 3.5e+62 years"; the other runs: 5.684 s and 6.954 s for 28 bits
\samedot
\begin{frame}{Task 3: what we saw}
  \begin{columns}[c,onlytextwidth]
    \column{0.3\textwidth}
      \small
      \renewcommand{\arraystretch}{1.4}
      \begin{tabular}{lr}
        16 bits & 0.002 s \\
        20 bits & 0.035 s \\
        24 bits & 0.34 s \\
        28 bits & 5.8 s \\
      \end{tabular}
    \column{0.67\textwidth}
      \centering
      \begin{tikzpicture}
        \begin{axis}[width=8.4cm, height=4.3cm, ybar, bar width=0.8cm, ymode=log, log origin=infty,
            ymin=0.001, ymax=60, symbolic x coords={16,20,24,28}, xtick=data, axis lines=left,
            axis line style={nLine}, tick label style={font=\scriptsize}, xlabel={\scriptsize bits of $p$},
            ylabel={\scriptsize seconds (log scale)}, ymajorgrids, grid style={nLight},
            nodes near coords, nodes near coords style={font=\scriptsize, text=nInk},
            point meta=explicit symbolic]
          \addplot[fill=DeepBlue!70, draw=DeepBlue] coordinates {(16,0.002)[0.002 s] (20,0.035)[0.035 s] (24,0.336)[0.34 s] (28,5.836)[5.8 s]};
        \end{axis}
      \end{tikzpicture}
  \end{columns}
  \stepbox{The last line}{\ttfamily 10374805 tries per second, a of 256 bits: about 3.5e+62 years}
  \curious{diffie-hellman/\#2.2}
  \keymsg{Every 4 more bits make the search about 16 times longer.}
  \say{Sprawdzone na tym laptopie (Intel Core i5-8265U), jedna seria pomiarów; studenci dostaną inne czasy. Znalezione a: 17773, 415826, 3465539, 60546499. Pętla: value = g, a = 1, dopóki value != A, value = value * g \% p i a += 1. Czasy zależą też od tego, jak duże było losowe a, dlatego wzrost nie jest równy; średnio około 16 razy na 4 bity. W symulatorze z linku te same cztery wyszukiwania w przeglądarce.}
\end{frame}

% CHECKED (python3): 60546499 / 5.836 = 10.4 million tries per second;
%   2**256 / 10374805 / (365*24*3600) = 3.5e62 years; 2**128 / 10374805 / (365*24*3600) = 1.0e24 years
\begin{frame}{From 28 bits to 256 bits}
  \begin{center}
    \begin{tikzpicture}
      \draw[nLine, very thick] (0,0) -- (6.2,0);
      \draw[nLine, very thick] (6.8,0) -- (12.5,0);
      \draw[nLine, very thick] (6.0,-0.25) -- (6.3,0.25) (6.6,-0.25) -- (6.9,0.25);
      \foreach \x/\b/\t in {0.4/16/0.002 s, 1.7/20/0.035 s, 3.0/24/0.34 s, 4.3/28/5.8 s} {
        \fill[DeepBlue] (\x,0) circle (0.12);
        \node[font=\scriptsize, text=nInk, above=3pt] at (\x,0.1) {\b\ bits};
        \node[font=\scriptsize, text=gray, below=3pt] at (\x,-0.1) {\t};
      }
      \node[font=\scriptsize, text=DeepBlue, align=center] at (2.35,1.0) {Task 3 on this laptop};
      \fill[AccentOrange] (12.2,0) circle (0.17);
      \node[font=\small\bfseries, text=AccentOrange, above=4pt] at (12.2,0.1) {$a$ of 256 bits};
      \node[font=\scriptsize, text=AccentOrange, below=4pt] at (12.2,-0.1) {$3.5 \times 10^{62}$ years};
    \end{tikzpicture}
  \end{center}
  \vspace{-0.1cm}
  {\small
  \begin{itemize}\setlength{\itemsep}{3pt}
    \item Our speed: about \textbf{10 million tries per second}. Every $a$ of 256 bits: $2^{256}$ tries, $3.5 \times 10^{62}$ years.
    \item Smarter general methods (about the square root of the number of possible $a$): about $2^{128}$ steps, still about $10^{24}$ years.
      That is why $a$ has 256 bits.
  \end{itemize}}
  \plainbox{Diffie-Hellman is safe as long as the \textbf{discrete logarithm is hard}.}
  \say{Wszechświat ma około 1,4 · 10\textasciicircum 10 lat. Metody ogólne (np. Pollard rho, baby-step giant-step) potrzebują około pierwiastka z liczby możliwych a, czyli 2\textasciicircum 128 kroków; krok jest droższy niż jedno mnożenie, więc 10\textasciicircum 24 lat to rząd wielkości. Te metody działają w każdej grupie, więc 256 bitów a daje 128 bitów bezpieczeństwa.}
\end{frame}

% CHECKED: NIST SP 800-57 Part 1 Rev. 5, table 2: FFC with p of 2048 bits = 112 bits of security, 3072 bits = 128,
%   7680 bits = 192, 15360 bits = 256; the same table gives RSA n of the same sizes (re-read 2026-11-25 in the PDF,
%   nvlpubs.nist.gov/nistpubs/SpecialPublications/NIST.SP.800-57pt1r5.pdf, table 2; 1024 bits: <= 80; FFC N = 224 at 112 bits)
\samedot
\begin{frame}{How big must p be?}
  \begin{center}\small The cleverest methods (the \textbf{number field sieve}) attack $p$ itself, not $a$.\\
    They are the \textbf{same family as for factoring}, so $p$ must be as big as RSA's $n$.\end{center}
  \vfill
  \begin{center}
    \renewcommand{\arraystretch}{1.5}
    \begin{tabular}{@{}lll@{}}
      \toprule
      \textbf{size of $p$} & \textbf{security} & \\
      \midrule
      1024 bits & about 80 bits & \textcolor{aWrite}{\faTimes\ too weak today} \\
      \rowcolor{AccentOrange!10} 2048 bits & \textbf{112 bits} & \textcolor{aExec}{\faCheck\ the minimum (RFC 3526 group 14)} \\
      3072 bits & 128 bits & \textcolor{aExec}{\faCheck\ as strong as a 128-bit key, or $a$ of 256 bits} \\
      \bottomrule
    \end{tabular}
  \end{center}
  \vfill
  \begin{center}\small ``Security of $n$ bits'': breaking it takes about $2^{n}$ steps.\end{center}
  \keymsg{$a$ needs 256 bits, but $p$ needs at least 2048: two different attacks set two sizes.}
  \say{Liczby z NIST SP 800-57 (część 1, tabela 2); te same rozmiary co dla n w RSA. Następny wykład: krzywe eliptyczne dają 128 bitów bezpieczeństwa przy 256 bitach zamiast 3072. Dla grupy 2048-bitowej 112 bitów to mniej niż 128 z wykładnika: najsłabsze ogniwo decyduje. Komputer kwantowy z algorytmem Shora łamałby logarytm dyskretny modulo p tak samo jak rozkład RSA, niezależnie od rozmiaru; o tym przy RSA i w następnym wykładzie.}
\end{frame}

\begin{frame}{When p is too small: Logjam (2015)}
  \vfill
  \begin{columns}[c,onlytextwidth]
    \column{0.3\textwidth}
      \begin{center}
      \begin{tikzpicture}
        \node[text=nInk, font=\fontsize{60}{60}\selectfont] at (0,0) {\faLock};
        \draw[white, line width=3pt] (-0.15,0.2) -- (0.1,-0.15) -- (-0.1,-0.45) -- (0.12,-0.9);
        \draw[aWrite, line width=1.4pt] (-0.15,0.2) -- (0.1,-0.15) -- (-0.1,-0.45) -- (0.12,-0.9);
        \node[draw=aWrite, very thick, fill=aWrite!8, rounded corners=3pt, font=\LARGE\bfseries, text=aWrite]
          at (0,-2.0) {512 bits};
      \end{tikzpicture}
      \end{center}
    \column{0.66\textwidth}
      Some servers still accepted old 512-bit \textbf{``export''} groups in TLS (the browser padlock, lecture 1):
      weakened on purpose in the 1990s, to obey US export rules.
      An attacker \textbf{on the connection} could make such a server pick one.
      \vspace{0.3cm}

      % CHECKED: Adrian et al., "Imperfect Forward Secrecy", CCS 2015 (weakdh.org/imperfect-forward-secrecy-ccs15.pdf), abstract:
      %   "After a week-long precomputation for a specified 512-bit group, we can compute arbitrary discrete logs in that group
      %   in about a minute"; section 1: "applies to 8.4% of Alexa Top Million HTTPS sites" (section 3: 8.4% supported DHE_EXPORT)
      %   (re-read 2026-11-25: pdftotext of the paper, lines quoted above found)
      \stepbox{The attack}{\faCalendar\ one week of computing for one fixed $p$,\\
        \faClock\ then about one minute for each discrete logarithm in that group}
      \vspace{0.2cm}

      The attack applied to \textbf{\textcolor{aWrite}{8.4\%}} of the top one million HTTPS web sites.
  \end{columns}
  \vfill
  \say{Grupy eksportowe to pozostałość po amerykańskich ograniczeniach z lat 90. Napastnik na połączeniu podmienia wiadomości na początku połączenia tak, że serwer wybiera grupę eksportową, a przeglądarka tego nie zauważa (downgrade); taką napastniczkę nazwiemy w sekcji 5 Mallory i tam do tego wrócimy. TLS to protokół zabezpieczający HTTPS. Atak opisali Adrian i inni na konferencji CCS w 2015 roku.}
\end{frame}

\samedot
\begin{frame}{One p for everybody}
  \begin{columns}[c,onlytextwidth]
    \column{0.5\textwidth}
      \begin{center}
      \begin{tikzpicture}[>=Stealth]
        \node[pub, font=\large\bfseries, minimum width=1.8cm, minimum height=1.0cm] (p) at (0,0) {one $p$};
        \foreach \i in {0,...,7} {
          \pgfmathsetmacro{\yy}{1.75-0.5*\i}
          \node[text=sTeal, font=\large] (s\i) at (3.6,\yy) {\faServer};
          \draw[data, thick] (p.east) -- (s\i.west);
        }
        \node[font=\scriptsize, text=nInk] at (3.6,-2.3) {millions of servers};
      \end{tikzpicture}
      \end{center}
    \column{0.48\textwidth}
      The week of work depends \textbf{only on $p$}, not on $a$ or $b$.
      \vspace{0.25cm}

      A few popular 1024-bit groups were used by millions of servers.
      \vspace{0.25cm}

      The Logjam authors estimated: a rich state could do the work \textbf{once} and read them all.
  \end{columns}
  \vfill
  \keymsg{At 2048 bits even the one-time work is out of reach, so one shared $p$ is fine. Use at least 2048 bits, or the elliptic version from the next lecture.}
  % CHECKED: Adrian et al., CCS 2015, abstract: "performing precomputation for a single 1024-bit group would allow passive
  %   eavesdropping on 18% of popular HTTPS sites, and a second group would allow decryption of traffic to 66% of IPsec VPNs
  %   and 26% of SSH servers"; section 1: ten groups: 66% of IKE VPNs, 26% of SSH, 16% of SMTP, 24% of popular HTTPS sites
  \say{Autorzy Logjam oszacowali: złamanie jednej, najpopularniejszej grupy 1024-bitowej pozwoliłoby podsłuchiwać 18\% popularnych stron HTTPS, a drugiej grupy: 66\% sieci VPN (IPsec) i 26\% serwerów SSH. Dlatego dziś p ma co najmniej 2048 bitów.}
\end{frame}

\begin{frame}{Good to know: the Diffie-Hellman problem}
  \vfill
  \begin{center}
    \begin{tikzpicture}[>=Stealth]
      \node[pub, font=\large] (ga) at (0,0.7) {$g^{a}$};
      \node[pub, font=\large] (gb) at (0,-0.7) {$g^{b}$};
      \node[hwmain, minimum width=2.4cm, minimum height=1.9cm, font=\Huge] (box) at (3.4,0) {?};
      \node[sec, font=\large] (gab) at (6.8,0) {$g^{ab}$};
      \draw[data] (ga.east) -- (box.west |- ga);
      \draw[data] (gb.east) -- (box.west |- gb);
      \draw[data] (box) -- (gab);
    \end{tikzpicture}
  \end{center}
  \vfill
  \begin{center}
    Eve does not really need $a$. She needs $g^{ab}$ from $g^{a}$ and $g^{b}$.\\[2pt]
    Nobody knows a way to do that faster than computing a discrete logarithm.
  \end{center}
  \vfill
  \say{Można pominąć. To formalnie inny problem (obliczeniowy problem Diffiego-Hellmana, CDH), ale w praktyce równie trudny. Gdyby ktoś pytał: dla szyfrowania ElGamal z sekcji 6 potrzeba jeszcze więcej, problemu decyzyjnego (DDH): Eve nie umie nawet odróżnić g\textasciicircum ab od losowej liczby. To założenie działa tylko w grupie o cyklu długości liczby pierwszej, jak w RFC 3526; w pełnym cyklu 22 Eve widzi parzystość, slajd ,,A cycle of 22 gives away one bit''.}
\end{frame}

% ======================================================================
\section[A key, not a cipher]{A key, not a cipher}
% ======================================================================

\begin{frame}{DH does not encrypt anything}
  \begin{center}\small The result of Diffie-Hellman is \textbf{one number} that both sides know.\end{center}
  \vfill
  \begin{center}
    \begin{tikzpicture}[>=Stealth]
      \node[hw, minimum width=3.8cm, minimum height=2.2cm] (dh) at (0,0) {};
      \node[text=nInk, font=\fontsize{30}{30}\selectfont] at (0,0.35) {\faHandshake};
      \node[font=\small\bfseries, text=nInk] at (0,-0.6) {Diffie-Hellman};
      \node[sec, font=\small] (s) at (3.6,0) {\faKey\ a shared secret};
      \node[hw, minimum width=3.8cm, minimum height=2.2cm] (c) at (7.6,0) {};
      \node[text=nInk, font=\fontsize{30}{30}\selectfont] at (7.6,0.35) {\faBolt};
      \node[font=\small\bfseries, text=nInk] at (7.6,-0.6) {ChaCha20 or AES};
      \draw[data] (dh) -- (s);
      \draw[data] (s) -- (c);
      \envpic{11.2}{0}{sBrown!15}{sBrown}
      \draw[data] (c) -- (10.35,0);
      \node[capt, below=0.6cm] at (11.2,0) {encrypts the messages with it};
      \node[capt, below=0.15cm of dh] {step 1};
      \node[capt, below=0.15cm of c] {step 2};
    \end{tikzpicture}
  \end{center}
  \vfill
  \keymsg{Diffie-Hellman agrees on a key. A symmetric cipher does the rest.}
  \say{Wymiana klucza to nie szyfrowanie: żadna wiadomość nie jest szyfrowana Diffie-Hellmanem.}
\end{frame}

% CHECKED (python3), reproducible: r = random.Random(2026); a = r.getrandbits(256); b = r.getrandbits(256);
%   secret = pow(pow(2, a, P), b, P) with P of RFC 3526 group 14: 616 digits, 726396308082692299...917295;
%   hashlib.sha256(secret.to_bytes(256, "big")).hexdigest() = f1584740e2d7aa28...8fac6398
\samedot
\begin{frame}{From a number to a key}
  \defbox{\textbf{Hash function:} any data in, a fixed-size fingerprint out, and it cannot be reversed. SHA-256 gives 32 bytes.}
  \begin{center}\small The secret has up to \textbf{2048 bits}. ChaCha20 needs \textbf{256}.\end{center}
  \vfill
  \begin{center}
    \begin{tikzpicture}[>=Stealth]
      \node[sec, font=\tiny\ttfamily, text width=3.6cm, align=left] (big) at (0,0)
        {\seqsplit{726396308082692299149293633956275253863230132969}\dots\seqsplit{917295}\\[2pt]\normalfont(616 digits)};
      % the funnel
      \draw[nInk, thick, fill=nFill] (2.7,1.0) -- (5.1,1.0) -- (4.3,-0.1) -- (4.3,-0.9) -- (3.5,-0.9) -- (3.5,-0.1) -- cycle;
      \node[font=\small\bfseries, text=nInk] at (3.9,0.45) {SHA-256};
      \draw[data] (big.east) -- (2.95,0.5);
      \node[sec, font=\tiny\ttfamily, text width=3.4cm, align=left] (key) at (8.0,-0.5)
        {\seqsplit{f1584740e2d7aa2868ceff74a6f36eea24507c49aafff59ade94c1a68fac6398}\\[2pt]\normalfont 32 bytes = 256 bits};
      \draw[data] (3.9,-0.9) |- (key.west);
    \end{tikzpicture}
  \end{center}
  \vfill
  \begin{center}
    \Large \texttt{key = SHA-256(secret)}\\[6pt]
    \small The secret is a number below $p$: its bits are \textbf{not evenly random}.\\
    The hash mixes all 2048 bits into 256 bits that look random.
  \end{center}
  \vfill
  \say{Liczba po lewej pochodzi z powtarzalnego przebiegu: a i b z random.Random(2026).getrandbits(256), grupa z RFC 3526; ten sekret ma 616 cyfr (sekret może mieć do 617 cyfr). Funkcja skrótu zamienia dane dowolnej długości w krótki odcisk. Sekret to liczba mniejsza od p, więc np. najwyższe bity nie są równomiernie losowe; klucz szyfru powinien mieć każdy bit losowy. Właściwe narzędzie do tego nazywa się KDF (funkcja wyprowadzania klucza) i będzie w następnym wykładzie. W zadaniu 4 robi to gotowa funkcja.}
\end{frame}

\begin{frame}{A new exchange for every session}
  \begin{center}\small \textbf{Ephemeral} Diffie-Hellman: $a$ and $b$ are new random numbers for every connection,
    and they are deleted when the connection ends.\end{center}
  \vfill
  \begin{center}
    \begin{tikzpicture}[>=Stealth]
      \foreach \i/\y in {1/0, 2/-1.2} {
        \node[text=DeepBlue, font=\Large] at (0,\y) {\faCalendar};
        \node[anchor=west, font=\small, text=nInk] at (0.5,\y) {connection \i:};
        \node[sec, font=\small] (s\i) at (4.0,\y) {$a_{\i},\ b_{\i}$};
        \node[sec, font=\small] (k\i) at (6.9,\y) {\faKey\ key \i};
        \draw[data] (s\i) -- (k\i);
        \draw[data, dashed] (k\i) -- (9.1,\y);
        \node[text=nInk, font=\Large] at (9.6,\y) {\faTrash};
        \node[anchor=west, font=\small, text=gray] at (10.1,\y) {(deleted)};
      }
      \node[font=\small, text=gray] at (4.0,-2.1) {\dots};
    \end{tikzpicture}
  \end{center}
  \vfill
  \say{Ephemeral znaczy ulotny. Liczenie nowej wymiany kosztuje mniej niż milisekundę, więc można to robić przy każdym połączeniu.}
\end{frame}

\samedot
\begin{frame}{Forward secrecy}
  \bigdef{\textbf{Forward secrecy:} a key stolen today does not open conversations recorded yesterday.}
  \vfill
  \begin{center}
    \begin{tikzpicture}[>=Stealth]
      % RSA: one stolen key opens every recording
      \node[anchor=west, font=\small\bfseries, text=nInk] at (-0.3,0.6) {RSA sends the key:};
      \node[sec, font=\small] (rk) at (0.6,-0.2) {\faKey\ stolen};
      \foreach \i in {0,1,2} {
        \node[ctx, font=\small] (r\i) at ({3.0+1.25*\i},-0.2) {\faLockOpen};
        \draw[data, thin] (rk.east) to[out=0, in=180] (r\i.west);
      }
      \node[anchor=west, font=\footnotesize, text=aWrite] at (6.6,-0.2) {steal the private key, decrypt every recorded session};
      % ephemeral DH
      \node[anchor=west, font=\small\bfseries, text=nInk] at (-0.3,-1.2) {ephemeral Diffie-Hellman:};
      \node[draw=nLight, dashed, thick, rounded corners=3pt, font=\small, text=gray] (ek) at (0.6,-2.0) {\faKey\ gone};
      \foreach \i in {0,1,2} \node[ctx, font=\small] at ({3.0+1.25*\i},-2.0) {\faLock};
      \node[anchor=west, font=\footnotesize, text=aExec] at (6.6,-2.0) {the session keys are gone, there is nothing to steal};
    \end{tikzpicture}
  \end{center}
  \vfill
  \begin{center}\small Eve's recording has only $A$ and $B$. Without $a$ or $b$ she needs a discrete logarithm,\\
    and $a$, $b$ no longer exist anywhere. Last lecture asked: \textit{what if a private key leaks?}\end{center}
  \keymsg{Ephemeral keys protect the past.}
  \say{To odpowiedź na problem z wykładu 3: wyciek klucza otwierał całą historię. Przy ulotnych kluczach tej historii nie ma czym otworzyć. Serwer ma też klucz długoterminowy (do pieczęci, sekcja 5). Jeśli go ukradną, Mallory może w przyszłości podszywać się pod serwer, ale nagranych rozmów nim nie otworzy: on tylko pieczętuje, nie szyfruje.}
\end{frame}

% CHECKED 2026-11-25: openssl speed -seconds 3 ffdh2048 ffdh3072, OpenSSL 3.0.13, the author's laptop (i5-8265U), three runs
%   (load average 0.2 to 2.2): 2048 bits 3801, 3849, 3920 op/s; 3072 bits 1519, 1484, 1520 op/s; 1/3850 = 0.26 ms, 1/1500 = 0.67 ms
%   (local/audyt/tmp-facts2/ossl.txt). An earlier run of 2026-11-19 gave 3713 and 1420.
\samedot
\begin{frame}{How fast is Diffie-Hellman?}
  \vfill
  \begin{center}
    \begin{tikzpicture}
      \begin{axis}[width=12cm, height=4.2cm, xbar, bar width=0.55cm, xmin=0, xmax=6000,
          symbolic y coords={3072-bit group,2048-bit group}, ytick=data, axis lines=left,
          axis line style={nLine}, tick label style={font=\small}, xlabel={\scriptsize operations per second (OpenSSL)},
          enlarge y limits=0.5, nodes near coords, point meta=explicit symbolic,
          nodes near coords style={font=\small, text=nInk, anchor=west}, nodes near coords align={horizontal}, xticklabel style={/pgf/number format/fixed}, scaled x ticks=false]
        \addplot[fill=DeepBlue!70, draw=DeepBlue] coordinates {(3800,2048-bit group)[about 3800 /s \ (0.26 ms)] (1500,3072-bit group)[about 1500 /s \ (0.67 ms)]};
      \end{axis}
    \end{tikzpicture}
  \end{center}
  \vfill
  \begin{center}Can we afford a new exchange for every connection? Yes: it costs \textbf{almost nothing}.\end{center}
  \vfill
  \say{Pomiar openssl speed ffdh2048 i ffdh3072 na tym laptopie z i5-8265U.}
\end{frame}

\begin{frame}{Where Diffie-Hellman is used}
  \vfill
  \begin{center}
    \begin{tikzpicture}
      \foreach \y/\ic/\name/\txt in {0/\faLock/TLS 1.3/every full connection setup includes an ephemeral exchange,
                                     -1.3/\faTerminal/SSH/the session key is agreed with Diffie-Hellman,
                                     -2.6/\faNetworkWired/IPsec/{IKE, the key exchange for many VPNs}} {
        \node[hw, minimum width=1.3cm, minimum height=1.0cm, font=\Large] at (0,\y) {\ic};
        \node[anchor=west, font=\large\bfseries, text=nInk] at (1.0,\y) {\name};
        \node[anchor=west, font=\normalsize, text=nInk] at (3.4,\y) {\txt};
      }
    \end{tikzpicture}
  \end{center}
  \vfill
  \begin{center}\small Today mostly in its \textbf{elliptic-curve} version: next lecture.\end{center}
  \vfill
  \say{SSH to szyfrowane logowanie na zdalny komputer, IPsec to szyfrowanie ruchu IP, na którym działa wiele sieci VPN. W TLS 1.3 każde pełne nawiązanie połączenia (full handshake) zawiera wymianę Diffiego-Hellmana; wyjątek to wznowienie sesji samym wcześniej ustalonym kluczem (tryb psk\_ke w RFC 8446, sekcja 4.2.9); jest też tryb psk\_dhe\_ke, czyli wznowienie z nową wymianą. Dziś coraz częściej wymiana jest hybrydowa: krzywa eliptyczna razem z algorytmem odpornym na komputery kwantowe (X25519MLKEM768), bo algorytm Shora złamałby sam Diffie-Hellman.}
\end{frame}

\begin{frame}{But who is on the other side?}
  \vfill
  \begin{center}
    \begin{tikzpicture}
      \alice{0}{0}
      \node[sec, font=\small] at (1.6,0.9) {\faKey};
      \draw[net] (Alice.east) -- (8.8,0);
      \node[circle, fill=nLight, text=white, minimum size=1.1cm, font=\Large] (q) at (9.5,0) {\faUser};
      \node[font=\Huge\bfseries, text=AccentOrange] at (10.6,0.3) {?};
    \end{tikzpicture}
  \end{center}
  \vfill
  \begin{center}
    Alice computed a secret \textbf{with somebody}.\\[4pt]
    \large How does she know it was Bob?
  \end{center}
  \vfill
  \say{Pytanie do sali przed następną sekcją.}
\end{frame}

% ======================================================================
\section[The man in the middle]{The man in the middle}
% ======================================================================

\begin{frame}{Mallory can change messages}
  \vfill
  \begin{center}
    \begin{tikzpicture}[>=Stealth]
      \alice{0}{0}
      \bob{12}{0}
      \draw[net] (Alice.east) -- (5.4,0);
      \draw[net] (6.6,0) -- (Bob.west);
      \mallory{6}{0}
      \node[text=sPinkText, font=\Large] at (6,1.0) {\faCut};
      \eve{3}{-1.7}
      \draw[data, dashed] (3,-0.15) -- (Eve.north);
      \node[text=sPinkText, font=\large, anchor=west] at (3.75,-1.7) {\faAssistiveListeningSystems};
    \end{tikzpicture}
  \end{center}
  \vfill
  \renewcommand{\arraystretch}{1.25}
  \begin{center}
    \begin{tabular}{ll}
      \textcolor{sPinkText}{\textbf{Eve}} & listens \\
      \textcolor{sPinkText}{\textbf{Mallory}} & sits on the cable: catches every message and sends what she wants \\
      & \small\textcolor{gray}{Like a router or a public Wi-Fi under her control.} \\
    \end{tabular}
  \end{center}
  \vfill
  \say{Eve jest bierna, Mallory aktywna. Taki atak nazywa się man-in-the-middle, człowiek pośrodku. Przypomnienie: Logjam z sekcji 3 to właśnie robota Mallory. Podmienia pierwsze wiadomości połączenia tak, że serwer wybiera grupę eksportową 512-bitową, a przeglądarka tego nie zauważa.}
\end{frame}

% the attack, one picture in three steps
\newcommand{\mitmscene}[1]{%
  \begin{tikzpicture}[>=Stealth]
    \useasboundingbox (-1,-2.95) rectangle (13,1.3);
    \alice{0}{0}
    \mallory{6}{0}
    \bob{12}{0}
    \draw[netarrow] (0.7,0.25) -- node[above, font=\small] {$A$} (5.3,0.25);
    \draw[netarrow] (6.7,0.25) -- node[above, font=\small, text=sPinkText] {$g^{m}$} (11.3,0.25);
    \node[sec, font=\scriptsize] at (6,1.05) {$m$: Mallory's secret};
    \ifnum#1>1
      \draw[netarrow] (11.3,-0.25) -- node[below, font=\small] {$B$} (6.7,-0.25);
      \draw[netarrow] (5.3,-0.25) -- node[below, font=\small, text=sPinkText] {$g^{m}$} (0.7,-0.25);
    \fi
    \ifnum#1>2
      \node[sec, font=\small] at (3,-1.6) {\faKey\ secret 1 $= g^{am}$};
      \node[sec, font=\small] at (9,-1.6) {\faKey\ secret 2 $= g^{bm}$};
      \node[font=\footnotesize, text=nInk, align=left, anchor=north] at (3,-2.0) {\textcolor{sIndigo}{\textbf{Alice:}} $(g^{m})^{a} = g^{am}$\\ \textcolor{sPinkText}{\textbf{Mallory:}} $A^{m} = g^{am}$};
      \node[font=\footnotesize, text=nInk, align=left, anchor=north] at (9,-2.0) {\textcolor{sTealText}{\textbf{Bob:}} $(g^{m})^{b} = g^{bm}$\\ \textcolor{sPinkText}{\textbf{Mallory:}} $B^{m} = g^{bm}$};
      \draw[decorate, decoration={brace, mirror, amplitude=5pt}, nInk] (0,-1.0) -- (6,-1.0);
      \draw[decorate, decoration={brace, mirror, amplitude=5pt}, nInk] (6,-1.0) -- (12,-1.0);
    \fi
  \end{tikzpicture}}
\newcommand{\mitmsteps}[1]{%
  \footnotesize
  \stepitem{1}{#1}{Alice sends $A$. Mallory keeps it and sends her own public value $g^{m}$ to Bob.}
  \stepitem{2}{#1}{Bob sends $B$. Mallory keeps it and sends $g^{m}$ to Alice.}
  \stepitem{3}{#1}{Alice shares a secret with Mallory. Bob shares another one with Mallory.}}

\begin{frame}{The attack, step 1}
  \vfill
  \begin{center}\mitmscene{1}\end{center}
  \vfill
  \begin{center}\begin{minipage}{0.8\textwidth}\mitmsteps{1}\end{minipage}\end{center}
  \vfill
  \say{Mallory ma własny sekret m i własną wartość publiczną g\textasciicircum m, tak jak Alice ma a i A. Zna m, więc policzy A\textasciicircum m: ma wspólny sekret z Alice.}
\end{frame}

\samedot
\begin{frame}{The attack, step 2}
  \vfill
  \begin{center}\mitmscene{2}\end{center}
  \vfill
  \begin{center}\begin{minipage}{0.8\textwidth}\mitmsteps{2}\end{minipage}\end{center}
  \vfill
  \say{Tak samo z drugiej strony: ma też wspólny sekret z Bobem.}
\end{frame}

\samedot
\begin{frame}{The attack, step 3}
  \vfill
  \begin{center}\mitmscene{3}\end{center}
  \vfill
  \begin{center}\begin{minipage}{0.8\textwidth}\mitmsteps{3}\end{minipage}\end{center}
  \vfill
  \say{Alice myśli, że g\textasciicircum m to B Boba, więc liczy (g\textasciicircum m)\textasciicircum a, tak jak na slajdzie ,,Why are they equal?''. Mallory dostaje tę samą liczbę jako A\textasciicircum m, bo zna m. Alice myśli, że jej sekret dzieli z Bobem. Bob myśli to samo. Oboje się mylą.}
\end{frame}

\samedot
\begin{frame}{What the attack gives}
  \vfill
  \begin{center}
    \begin{tikzpicture}[>=Stealth]
      \alice{0}{0}
      \envpic{2.7}{0}{sBrown!15}{sBrown}
      \node[capt, below=0.55cm] at (2.7,0) {encrypted with key 1};
      \draw[data] (Alice) -- (1.9,0);
      \draw[data] (3.5,0) -- (5.0,0);
      \mallory{6}{0}
      \node[draw=sPink, thick, fill=white, rounded corners=3pt, font=\scriptsize, text=nInk, align=left, anchor=south]
        at (6,0.95) {\faLockOpen\ decrypts\\ \faEye\ reads\\ \faLock\ encrypts again};
      \draw[data] (7.0,0) -- (8.5,0);
      \envpic{9.3}{0}{sBrown!15}{sBrown}
      \node[capt, below=0.55cm] at (9.3,0) {encrypted with key 2};
      \draw[data] (10.1,0) -- (11.4,0);
      \bob{12.1}{0}
      \node[text=aExec, font=\large] at (12.1,1.0) {\faCheck};
    \end{tikzpicture}
  \end{center}
  \vfill
  \begin{center}Both sides see \textbf{correct messages}. Nothing looks wrong.\end{center}
  \vfill
  \curious{diffie-hellman/\#3}
  \keymsg{Diffie-Hellman alone cannot tell who is on the other side.}
  \say{Szyfrowanie działa bez zarzutu, tylko z niewłaściwą osobą. W symulatorze z linku atak krok po kroku, na tych samych małych liczbach.}
\end{frame}

\begin{frame}{Task 4: be Mallory}
  \vfill
  \taskgoal{sit between Alice and Bob, read Alice's message, and let Bob notice nothing.}
  \taskready{\texttt{network.py}: Alice, Bob, \texttt{key\_from\_secret}, \texttt{encrypt} and \texttt{decrypt}.}
  \vfill
  \opentask{Task4}{about 10 minutes}
  \vfill
  \say{Około 10 minut. network.py jest gotowy i nie trzeba go zmieniać: zawiera Alice, Boba i szyfrowanie wiadomości kluczem z SHA-256. Bez zmian program wypisuje tylko to, co czyta Bob.}
\end{frame}

% CHECKED: Task4/solution/mitm.py prints the two lines below
\samedot
\begin{frame}{Task 4: what we saw}
  \vfill
  \begin{center}
    \begin{minipage}{0.86\textwidth}
      \stepbox{The output}{\ttfamily \textcolor{sPinkText}{Mallory reads:} My new card PIN is 7314. Please keep it secret.\\
        \textcolor{sTealText}{Bob reads:}\ \ \ \ \,My new card PIN is 7314. Please keep it secret.}
    \end{minipage}\\[0.4cm]
    Mallory read the PIN. Bob got the right message, so nothing looked wrong.
  \end{center}
  \vfill
  \keymsg{Two exchanges, two keys, and Mallory holds both.}
  \say{Sprawdzone. change\_A: self.key\_alice = key\_from\_secret(pow(A, self.m, P)) i zwraca self.g\_m (wartość publiczną Mallory, g\textasciicircum m) zamiast A. change\_B tak samo, z kluczem dla Boba. change\_message: decrypt kluczem Alice, print, encrypt kluczem Boba. Bob dostał poprawną wiadomość i niczego nie zauważył.}
\end{frame}

% Mallory replaces both public values by one bad value: \badscene{value}{Alice's computation}{Bob's computation}{Mallory's box}
\newcommand{\badscene}[4]{%
  \begin{tikzpicture}[>=Stealth]
    \useasboundingbox (-1,-2.7) rectangle (13,1.4);
    \alice{0}{0}
    \mallory{6}{0}
    \bob{12}{0}
    \draw[netarrow] (0.7,0.25) -- node[above, font=\small] {$A = 8$} (5.3,0.25);
    \draw[netarrow] (6.7,0.25) -- node[above, font=\small\bfseries, text=sPinkText] {#1} (11.3,0.25);
    \draw[netarrow] (11.3,-0.25) -- node[below, font=\small] {$B = 19$} (6.7,-0.25);
    \draw[netarrow] (5.3,-0.25) -- node[below, font=\small\bfseries, text=sPinkText] {#1} (0.7,-0.25);
    \node[draw=sPink, thick, fill=white, rounded corners=3pt, font=\scriptsize, text=nInk, align=center, anchor=south] at (6,0.75) {#4};
    \node[sec, font=\small] at (1.2,-1.9) {#2};
    \node[sec, font=\small] at (10.8,-1.9) {#3};
    \node[capt] at (1.2,-2.5) {\textcolor{sIndigo}{\textbf{Alice}}, $a = 6$};
    \node[capt] at (10.8,-2.5) {\textcolor{sTealText}{\textbf{Bob}}, $b = 15$};
  \end{tikzpicture}}

% CHECKED (python3): pow(1, 6, 23) = 1, pow(1, 15, 23) = 1
\begin{frame}{Mallory sends 1}
  \vfill
  \begin{center}\badscene{1}{$1^{6} \bmod 23 = \mathbf{1}$}{$1^{15} \bmod 23 = \mathbf{1}$}{no secret $m$ needed:\\ $1^{\text{anything}} = 1$}\end{center}
  \vfill
  \begin{center}
    Alice and Bob get \textbf{the same} secret, 1, so they even have the same key.\\[2pt]
    Mallory knows it without computing anything. She reads every message and passes it on unchanged.
  \end{center}
  \vfill
  \keymsg{A value from the other side can be a trap. Even ``we have the same key'' proves nothing.}
  \say{Na tych samych liczbach co w wymianie: A = 8 i B = 19, ale Mallory obie podmienia na 1. W zadaniu 4 Mallory musiała prowadzić dwie wymiany i szyfrować na nowo; tu nie liczy nic. Gdyby Alice i Bob porównali swoje klucze (np. odcisk klucza przez telefon), zobaczyliby to samo i uspokoiliby się. Pytanie do sali: jakie jeszcze liczby byłyby tak samo złe? Następny slajd.}
\end{frame}

% CHECKED (python3): pow(22, 6, 23) = 1, pow(22, 15, 23) = 22
\samedot
\begin{frame}{Mallory sends p - 1}
  \vfill
  \begin{center}\badscene{22}{$22^{6} \bmod 23 = \mathbf{1}$}{$22^{15} \bmod 23 = \mathbf{22}$}{the secret is 1 or 22:\\ two tries}\end{center}
  \vfill
  \begin{center}
    Recall ``A bad g'': the powers of $22 = p - 1$ are only 22 and 1.\\[2pt]
    The same bad number, but now \textbf{Mallory} picks it, not the standard.
  \end{center}
  \vfill
  \keymsg{A value with a short cycle traps the secret among a few values.}
  \say{Alice ma a = 6, parzyste, więc dostaje 1; Bob ma b = 15, nieparzyste, więc 22. Mallory nie zna a ani b, ale ma tylko dwa kandydaty na każdy sekret i próbuje oba. Ogólnie: wartość o krótkim cyklu zamyka sekret w małym zbiorze (small-subgroup confinement, zamknięcie w małej podgrupie).}
\end{frame}

% CHECKED: RFC 8446 (TLS 1.3), section 4.2.8.1: "Peers MUST validate each other's public key Y by ensuring that 1 < Y < p-1.
%   This check ensures that the remote peer is properly behaved and isn't forcing the local system into a small subgroup."
%   RFC 7919, section 5.1: the same range check; NIST SP 800-56A Rev. 3, 5.6.2.3.1 (full validation): 2 <= y <= p - 2 and y^q mod p = 1
% CHECKED (python3): p = 23, q = 11, g = 2: pow(16, 11, 23) = 1 (16 = 2^4), pow(22, 11, 23) = 22, pow(5, 11, 23) = 22
\samedot
\begin{frame}{Check what you receive}
  \vfill
  \begin{center}
    \begin{tikzpicture}
      \node[pub, font=\Large, minimum width=2.0cm, minimum height=1.1cm] (v) at (0,0) {$B$};
      \node[draw=aExec, very thick, fill=aExec!8, rounded corners=4pt, font=\Large, inner sep=8pt] (c) at (4.6,0) {$1 < B < p - 1$};
      \draw[data] (v) -- (c);
      \node[text=aExec, font=\Large, anchor=west] (ok) at (7.4,0.45) {\faCheck};
      \node[font=\small, text=nInk, anchor=west] at (8.0,0.45) {go on};
      \node[text=aWrite, font=\Large, anchor=west] at (7.4,-0.45) {\faTimes};
      \node[font=\small, text=nInk, anchor=west] at (8.0,-0.45) {1 or $p - 1$: stop};
    \end{tikzpicture}
  \end{center}
  \vfill
  \begin{itemize}\setlength{\itemsep}{5pt}
    \item TLS 1.3 requires it from both sides: ``Peers MUST validate each other's public key'' (RFC 8446).
    \item With a safe prime, 1 and $p - 1$ are the only values with a short cycle, so this check is enough.
    \item The full check of NIST also tests $B^{q} \bmod p = 1$: $B$ lies on the cycle of $g$.\\
      \small Recall ``Why a cycle of prime length?'': with $g = 2$ mod 23, $q = 11$, and $16^{11} \bmod 23 = 1$, but $22^{11} \bmod 23 = 22$.
  \end{itemize}
  \vfill
  \keymsg{Never use a public value before you check it.}
  \say{Dlaczego wystarcza przy bezpiecznej liczbie pierwszej p = 2q + 1: długość cyklu każdej liczby dzieli p - 1 = 2q, więc może być 1 (tylko liczba 1), 2 (tylko p - 1), q albo 2q. Liczby z cyklem q lub 2q są długie, więc wystarczy odrzucić 1 i p - 1. Uwaga: sprawdzenie B do potęgi q ma sens tylko w grupie z cyklem q, jak w RFC 3526; przy g = 5 mod 23 (pełny cykl) prawdziwe B = 19 daje 19\textasciicircum 11 mod 23 = 22 i zostałoby odrzucone. Sprawdzenie nie zatrzyma Mallory z zadania 4: jej g\textasciicircum m jest porządną liczbą. Na to potrzebna jest pieczęć, za chwilę.}
\end{frame}

\begin{frame}{Task 5: Mallory sends 1}
  \vfill
  \taskgoal{read Alice's message with no secret of your own, then stop the attack with one check.}
  \taskready{\texttt{network.py}: Alice and Bob, who call your check on every value they receive.}
  \begin{center}
    \small\textcolor{AccentOrange}{\faLightbulb}\ \textbf{Predict first:} must Mallory encrypt the message again for Bob, as in Task 4?
  \end{center}
  \vfill
  \opentask{Task5}{about 8 minutes}
  \vfill
  \say{Około 8 minut. Program uruchamia wymianę dwa razy: bez sprawdzania i z funkcją check\_public studenta. Zapisać przewidywania na tablicy.}
\end{frame}

% CHECKED: python3 Task5/solution/send_one.py (PYTHONPATH=Task5), 2026-11-19, prints the four lines below
\samedot
\begin{frame}{Task 5: what we saw}
  \vfill
  \begin{center}
    \begin{minipage}{0.86\textwidth}
      \stepbox{The output}{\ttfamily Without the check:\\
        \textcolor{sPinkText}{Mallory reads:} The new Wi-Fi password is Blue-Heron-58.\\
        \textcolor{sTealText}{Bob reads:}\ \ \ \ \,The new Wi-Fi password is Blue-Heron-58.\\
        With the check:\\
        \textcolor{aWrite}{Bob rejects A = 1 and stops.}}
    \end{minipage}\\[0.4cm]
    Both keys came from the secret 1, so Mallory did not even re-encrypt.\\
    One comparison, $1 < A < p - 1$, stopped the attack before Bob computed anything.
  \end{center}
  \vfill
  \curious{diffie-hellman/\#3}
  \keymsg{Without a check, Mallory needs no secret at all.}
  \say{Sprawdzone. change\_A i change\_B zwracają 1, change\_message: self.key = key\_from\_secret(1), print, i zwraca dane bez zmian. check\_public: return 1 < value < P - 1. W symulatorze z linku przełącznik ,,Mallory sends: 1 / p - 1'' i sprawdzanie 1 < value < p - 1, na małych liczbach.}
\end{frame}

% CHECKED: C. H. Lim, P. J. Lee, "A key recovery attack on discrete log-based schemes using a prime order subgroup",
%   CRYPTO '97; CVE-2016-0701, OpenSSL advisory 2016-01-28 (openssl-library.org/news/secadv/20160128.txt): OpenSSL 1.0.2,
%   X9.42 parameters as for RFC 5114 with primes that are not safe, the server reusing its private DH exponent;
%   fixed in 1.0.2f; reported by Antonio Sanso (Adobe). L. Valenta et al., "Measuring small subgroup attacks against
%   Diffie-Hellman", NDSS 2017
\samedot
\begin{frame}{Good to know: when the secret stays the same}
  \vfill
  \begin{columns}[c,onlytextwidth]
    \column{0.36\textwidth}
      \begin{center}
      \begin{tikzpicture}
        \node[text=sTeal, font=\fontsize{44}{44}\selectfont] at (0,0.5) {\faServer};
        \node[sec, font=\small] at (0,-0.85) {the same $b$\\ for every client};
        \foreach \i/\y in {1/1.5, 2/0.6, 3/-0.3} {
          \node[text=sPink, font=\large] (m\i) at (-2.4,\y) {\faUserNinja};
          \draw[data, thin] (m\i) -- (-0.75,0.5);
        }
        \node[font=\scriptsize, text=sPinkText, align=center] at (-2.4,-0.95) {many\\ handshakes};
      \end{tikzpicture}
      \end{center}
    \column{0.6\textwidth}
      \small
      A server that keeps \textbf{one} secret $b$ for many connections (a \textbf{static} key, not ephemeral):
      \vspace{0.2cm}

      \begin{itemize}\setlength{\itemsep}{4pt}
        \item If $p - 1$ has many small factors, Mallory sends values with short cycles of length 2, 3, 5, \dots{} and learns $b$ modulo each of them, a few handshakes each (Lim and Lee, 1997).
        \item Put the pieces together, and she has the server's \textbf{whole} $b$.
        \item \textbf{CVE-2016-0701}: OpenSSL 1.0.2 with the groups of RFC 5114 (not safe primes) and a reused $b$. Fixed in January 2016.
      \end{itemize}
  \end{columns}
  \vfill
  \keymsg{Safe prime, checked values, a new secret every time: each closes the same door.}
  \say{Można pominąć. Klucz statyczny to odwrotność ulotnego: serwer mógł trzymać jedno b dla wielu połączeń (OpenSSL 1.0.2 robił tak, gdy opcja SSL\_OP\_SINGLE\_DH\_USE była wyłączona), żeby oszczędzić czas. Wtedy każda zła wartość od Mallory zdradza kawałek b, jak parzystość w slajdzie ,,A cycle of 22 gives away one bit''. Błąd zgłosił Antonio Sanso (Adobe); poprawka w OpenSSL 1.0.2f sprawdza B\textasciicircum q i zawsze bierze nowe b. Valenta i inni (NDSS 2017): z ponad 20 bibliotek żadna przed styczniem 2016 nie sprawdzała domyślnie rzędu podgrupy.}
\end{frame}

% a seal: \seal{x}{y}{colour}{text}
\newcommand{\seal}[4]{%
  \node[circle, draw=#3, very thick, fill=#3!15, text=#3, minimum size=1.0cm, inner sep=0pt, font=\tiny\bfseries, align=center]
    at (#1,#2) {\faStamp\\#4};}
\begin{frame}{The fix: authenticate the exchange}
  \vfill
  \begin{center}
    \begin{tikzpicture}
      \node[pub, font=\Large, minimum width=2.2cm, minimum height=1.3cm] at (0,0) {$B$};
      \seal{1.75}{0}{sTeal}{Bob}
      \node[text=aExec, font=\Large] at (0,-1.4) {\faCheck};
      \node[capt] at (0,-1.95) {sealed by Bob};
      \node[pub, font=\Large, minimum width=2.2cm, minimum height=1.3cm, text=sPinkText] at (5,0) {$g^{m}$};
      \node[circle, draw=nLight, very thick, dashed, minimum size=1.0cm] at (6.75,0) {};
      \node[text=aWrite, font=\Large] at (5,-1.4) {\faTimes};
      \node[capt] at (5,-1.95) {Mallory has no seal of Bob};
    \end{tikzpicture}
  \end{center}
  \vfill
  \begin{center}
    Bob adds to his $B$ a \textbf{seal} that only his long-term private key (not $b$) can make.\\
    Alice checks the seal with Bob's public key.\\[6pt]
    Mallory can replace $B$, but she \textbf{cannot seal} her $g^{m}$ as Bob.\\[4pt]
    \small The seal does not replace the check $1 < B < p - 1$: TLS 1.3 requires both.\\[2pt]
    \small\textcolor{gray}{Such a seal is a digital signature: lecture 8.}
  \end{center}
  \vfill
  \curious{diffie-hellman/\#3}
  \say{Tylko idea: pieczęć robi się długoterminowym kluczem prywatnym Boba (tym samym od lat, nie ulotnym b), sprawdza kluczem publicznym, a Mallory nie ma klucza prywatnego Boba. W symulatorze z linku można włączyć sprawdzanie pieczęci i zobaczyć, że atak się zatrzymuje. Jak dokładnie działają podpisy cyfrowe, to osobny temat, którego tu nie omawiamy.}
\end{frame}

\samedot
\begin{frame}{How does Alice know Bob's public key?}
  \begin{columns}[c,onlytextwidth]
    \column{0.42\textwidth}
      \begin{center}
      \begin{tikzpicture}
        \draw[nLine, thick, fill=white, drop shadow={opacity=0.25}] (0,0) rectangle (4.6,3.2);
        \draw[DeepBlue, thick] (0.15,0.15) rectangle (4.45,3.05);
        \node[font=\small\bfseries, text=DeepBlue] at (2.3,2.65) {Bob's public key};
        \node[font=\scriptsize, text=nInk, anchor=west] at (0.4,2.1) {name: \textbf{Bob}};
        \node[font=\scriptsize, text=nInk, anchor=west] at (0.4,1.65) {key: \textbf{Bob's public key}};
        \node[circle, draw=AccentOrange, very thick, fill=AccentOrange!15, text=AccentOrange, minimum size=1.2cm,
          font=\tiny\bfseries, align=center] at (3.6,0.85) {\faStamp\\trusted};
        \node[font=\tiny, text=gray, anchor=west] at (0.4,0.6) {vouched for by a trusted party};
      \end{tikzpicture}
      \end{center}
    \column{0.55\textwidth}
      Bob's public key comes \textbf{vouched for} by someone Alice already trusts.
      \vspace{0.4cm}

      \textcolor{DeepBlue}{\faGlobe}\ Browsers come with a list of trusted parties.
      \vspace{0.4cm}

      \small\textcolor{gray}{How this works: later in the course.}
  \end{columns}
  \vfill
  \keymsg{Diffie-Hellman needs authentication. Without it, it talks to anyone.}
  \say{Tylko pojęcie, bez szczegółów. Takie poręczenie to certyfikat; dokładnie omawiają to wykłady 9 do 12 (PKI), więc tu nie używamy tego słowa jako znanego. Dlatego przeglądarka ostrzega, gdy poręczenie strony się nie zgadza.}
\end{frame}

\samedot
\begin{frame}{In TLS 1.3}
  \vfill
  \begin{center}
    \begin{tikzpicture}
      \foreach \y/\n/\ic/\a/\b in {1.3/1/\faHandshake/Ephemeral Diffie-Hellman/a fresh shared key,
                                0/2/\faStamp/The server's seal/proves who sent the server's value,
                                -1.3/3/\faCertificate/A trusted party vouches/for whose public key that is} {
        \node[hw, minimum width=5.6cm, minimum height=1.05cm, anchor=west, align=left] at (0,\y) {};
        \node[anchor=west] at (0.15,\y) {\bigbadge{\n}};
        \node[anchor=west, font=\large, text=nInk] at (0.8,\y) {\ic};
        \node[anchor=west, font=\small\bfseries, text=nInk] at (1.5,\y) {\a};
        \draw[data] (5.8,\y) -- (6.8,\y);
        \node[anchor=west, font=\small, text=nInk] at (6.9,\y) {\b};
      }
      \draw[container] (-0.2,-2.0) rectangle (12.6,2.0);
      \node[fill=white, font=\scriptsize\bfseries, text=nLine] at (1.3,2.0) {one protocol};
    \end{tikzpicture}
  \end{center}
  \vfill
  \say{To dokładnie ten przepis: wymiana daje tajność, pieczęć serwera i poręczenie jego klucza (certyfikat, wykłady 9 do 12) dają pewność, z kim rozmawiamy.}
\end{frame}

% ======================================================================
\section[ElGamal encryption]{ElGamal encryption}
% ======================================================================

\begin{frame}{Diffie-Hellman as encryption}
  \begin{columns}[c,onlytextwidth]
    \column{0.52\textwidth}
      \textbf{Taher ElGamal, 1985:}
      \vspace{0.2cm}

      \begin{itemize}\setlength{\itemsep}{6pt}
        \item Bob publishes his half of the exchange, $y = g^{x}$, \textbf{once and for all}.
        \item For every message, Alice does her half \textbf{on the spot}
          and uses the shared value to hide the message.
        \item Why: Bob may be offline (e-mail). Alice cannot wait for his $B$, so Bob publishes it in advance.
      \end{itemize}
    \column{0.46\textwidth}
      \begin{center}
      \begin{tikzpicture}[>=Stealth]
        % the notice board
        \draw[nInk, very thick, fill=nFill, rounded corners=2pt] (0,0) rectangle (3.4,2.4);
        \draw[nInk, very thick] (0.4,0) -- (0.4,-0.6) (3.0,0) -- (3.0,-0.6);
        \node[pub, font=\large, rotate=-4] (y) at (1.7,1.25) {$y = g^{x}$};
        \fill[aWrite] (1.7,1.75) circle (0.08);
        \node[font=\tiny\bfseries, text=nInk] at (1.7,2.15) {notice board};
        \bob{4.4}{1.3}
        \alice{-0.2}{-1.6}
        \node[msgb, font=\small] (m) at (1.9,-1.6) {\faEnvelope\ message};
        \draw[data] (m.east) to[out=0, in=-90] (Bobname.south);
      \end{tikzpicture}
      \end{center}
  \end{columns}
  \say{ElGamal to Diffie-Hellman, w którym jedna strona jest stała, a druga robi swoją połowę razem z wiadomością.}
\end{frame}

% CHECKED (python3): pow(5, 15, 23) = 19
\begin{frame}{The keys}
  \vfill
  \begin{center}
    \begin{tikzpicture}
      \node[text=sAmber, font=\fontsize{40}{40}\selectfont] at (0,0) {\faKey};
      \node[font=\small\bfseries, text=sAmberText] at (0,-1.0) {private key};
      \node[text=nLine, font=\fontsize{40}{40}\selectfont] at (5,0) {\faKey};
      \node[font=\small\bfseries, text=nInk] at (5,-1.0) {public key};
    \end{tikzpicture}
  \end{center}
  \vfill
  \begin{center}
    \renewcommand{\arraystretch}{1.7}
    \begin{tabular}{lll}
      \textcolor{sAmberText}{\textbf{private key}} & $x$ & \tikz[baseline=(n.base)]\node[sec] (n) {$x = 15$}; \\
      \textbf{public key} & $y = g^{x} \bmod p$ & \tikz[baseline=(n.base)]\node[pub] (n) {$y = 5^{15} \bmod 23 = 19$}; \\
      public too & $p$ and $g$ & \tikz[baseline=(n.base)]\node[pub] (n) {$p = 23,\ g = 5$}; \\
    \end{tabular}
  \end{center}
  \vfill
  \say{Ten sam przykład co wcześniej: x = 15 to b Boba z wymiany, a y = 19 to jego B. Bob naprawdę publikuje swoją połowę wymiany.}
\end{frame}

% CHECKED (python3): pow(5, 3, 23) = 10, pow(19, 3, 23) = 5, 7 * 5 % 23 = 12
\samedot
\begin{frame}{Encrypt}
  \vfill
  \begin{center}
    \renewcommand{\arraystretch}{1.8}
    \begin{tabular}{l@{\hspace{1.2cm}}l}
      1. pick a random secret $k$ & \tikz[baseline=(n.base)]\node[sec] (n) {$k = 3$}; \\
      2. $C_1 = g^{k} \bmod p$ & $C_1 = 5^{3} \bmod 23 = 10$ \\
      3. $C_2 = M \times y^{k} \bmod p$ & \tikz[baseline=(n.base)]\node[msgb, inner sep=2pt] (n) {$M = 7$};:\quad
        $C_2 = 7 \times 19^{3} \bmod 23 = 12$ \\
    \end{tabular}
  \end{center}
  \vfill
  \begin{center}
    The ciphertext is the pair\quad \tikz[baseline=(n.base)]\node[ctx, font=\Large] (n) {$(C_1, C_2) = (10, 12)$};\\[0.3cm]
    \footnotesize\textcolor{gray}{$M$ is the message as a number smaller than $p$: text $\to$ bytes $\to$ one big number, as in RSA.}
  \end{center}
  \vfill
  \curious{diffie-hellman/\#4}
  \say{19\textasciicircum 3 = 6859, a 6859 mod 23 = 5, więc C2 = 7 · 5 mod 23 = 35 mod 23 = 12. Wiadomość M to liczba mniejsza od p. W symulatorze z linku ten sam przykład: szyfrowanie i odszyfrowanie. Karta ,,ta sama wiadomość dwa razy'' jest tam zwinięta pod przyciskiem; nie otwierać jej przed zadaniem 6.}
\end{frame}

\samedot
\begin{frame}{Recall: the modular inverse}
  \vfill
  \begin{center}
    the inverse of $s \bmod p$:\quad the number $s^{-1}$ with\quad $s \times s^{-1} \bmod p = 1$
  \end{center}
  \vfill
  \begin{center}
    \Large the inverse of $5 \bmod 23$ is $\mathbf{14}$,\quad because\quad
    $5 \times 14 = 70 = 3 \times 23 + \colorbox{AccentOrange!25}{$\mathbf{1}$}$
  \end{center}
  \vfill
  \begin{center}
    \normalsize In Python (3.8 or newer):\quad \pyprompt{pow(5, -1, 23)}{14}
  \end{center}
  \vfill
  % CHECKED (python3): pow(5, -1, 23) = 14
  \curious{math/extended-euclid/\#3}
  \say{Odwrotność znamy z RSA: tam d = pow(e, -1, phi). Działa jak dzielenie modulo p. Python liczy ją rozszerzonym algorytmem Euklidesa; jak to działa, pokazuje strona z linku.}
\end{frame}

% CHECKED (python3): pow(10, 15, 23) = 5, pow(5, -1, 23) = 14, 12 * 14 % 23 = 7
\samedot
\begin{frame}{Decrypt}
  \vfill
  \begin{center}
    \renewcommand{\arraystretch}{1.8}
    \begin{tabular}{l@{\hspace{1.2cm}}l}
      1. $s = C_1^{\,x} \bmod p$ & $s = 10^{15} \bmod 23 = 5$ \\
      2. $s^{-1} \bmod p$ & $5^{-1} \bmod 23 = 14$ \\
      3. $M = C_2 \times s^{-1} \bmod p$ & $M = 12 \times 14 \bmod 23 = 168 \bmod 23 = 7$ \\
    \end{tabular}
  \end{center}
  \vfill
  \begin{center}
    The message is back:\quad \tikz[baseline=(n.base)]\node[msgb, font=\Large] (n) {$M = 7$};
  \end{center}
  \vfill
  \say{Wszystkie liczby sprawdzone w Pythonie.}
\end{frame}

% CHECKED (python3): pow(10, 15, 23) = 5 and pow(19, 3, 23) = 5
\samedot
\begin{frame}{Why does it work?}
  \begin{center}
    \begin{tikzpicture}
      \envpic{0}{0}{sBrown!15}{sBrown}
      \node[font=\large\bfseries, text=nInk, fill=sBrown!15, inner sep=1pt] at (0,-0.28) {$C_1$};
      \node[capt, anchor=east, align=right] at (-1.0,0) {$C_1 = g^{k}$: Alice's half of\\ a Diffie-Hellman exchange};
      \envpic{2.0}{0}{sBrown!15}{sBrown}
      \node[font=\large\bfseries, text=nInk, fill=sBrown!15, inner sep=1pt] at (2.0,-0.28) {$C_2$};
      \node[capt, anchor=west, align=left] at (3.0,0) {$C_2 = M \times y^{k}$: the message,\\ hidden by the shared secret};
      \draw[nInk, line width=2pt, rounded corners=2pt] (0.75,0.65) -- (0.75,0.25) -- (1.25,0.25) -- (1.25,0.65);
    \end{tikzpicture}
  \end{center}
  \vspace{-0.1cm}
  \begin{center}
    \begin{tikzpicture}
      \node[anchor=east, font=\large\bfseries, text=sTealText] at (0,0.6) {Bob computes};
      \node[anchor=west, font=\Large] at (0.2,0.6) {$C_1^{\,x} \;=\; (g^{k})^{x} \;=\; g^{\,kx}$};
      \node[anchor=east, font=\large\bfseries, text=sIndigo] at (0,-0.6) {Alice used};
      \node[anchor=west, font=\Large] at (0.2,-0.6) {$y^{k} \;=\; (g^{x})^{k} \;=\; g^{\,xk}$};
      \node[sec, font=\Large, minimum height=2.1cm, minimum width=2.6cm] (r) at (8.0,0) {$g^{kx} \bmod p$};
      \draw[data] (5.0,0.6) -- (r.west |- 0,0.6);
      \draw[data] (5.0,-0.6) -- (r.west |- 0,-0.6);
    \end{tikzpicture}
  \end{center}
  \begin{center}\small In our example: $10^{15} \bmod 23 = 5$ and $19^{3} \bmod 23 = 5$. Without $C_1$, Bob could not find it.\\[2pt]
    $y^{k}$ looks random to Eve (she would need $k$ or $x$), so $M \times y^{k}$ looks random, like XOR with a keystream.\end{center}
  \keymsg{Alice multiplies $M$ by the shared secret. Bob divides by it.}
  \say{Dzielenie modulo p to mnożenie przez odwrotność. Dwie części szyfrogramu: C1 to połowa wymiany od Alice (Bob potrzebuje jej, żeby policzyć sekret), C2 to wiadomość zakryta tym sekretem. Uczciwie: ,,wygląda losowo'' jest prawdą w grupie o cyklu długości liczby pierwszej, jak w RFC 3526 i w zadaniu 6. Przy g = 5 mod 23 (pełny cykl 22) Eve widzi z C1 i y, czy y\textasciicircum k jest ,,kwadratem'', i z tego jeden bit o M, jak parzystość w sekcji 3.}
\end{frame}

\begin{frame}{Task 6: ElGamal in Python}
  \vfill
  \taskgoal{write \texttt{encrypt} and \texttt{decrypt} of ElGamal, for the real 2048-bit group.}
  \taskready{Bob's keys $x$ and $y$, and text to number and back, in \texttt{elgamal.py}.}
  \vfill
  \opentask{Task6}{about 8 minutes}
  \vfill
  \say{Około 8 minut. Wzory są na poprzednich slajdach. Tekst zamieniany na liczbę i z powrotem jest już gotowy w pliku. Odwrotność: pow(s, -1, P).}
\end{frame}

% CHECKED: Task6/solution/elgamal.py prints the two lines below (and C1, C2: new digits on every run)
\samedot
\begin{frame}{Task 6: what we saw}
  \vfill
  \begin{center}
    \begin{minipage}{0.8\textwidth}
      \stepbox{The output}{\ttfamily decrypted: Meet me at the old bridge at noon.\\
        the same message again gives the same C1: False}
    \end{minipage}\\[0.4cm]
    The message comes back. Encrypted twice, it gives two different ciphertexts.
  \end{center}
  \vfill
  \keymsg{Every call picks a new $k$, so every ciphertext is new.}
  \say{Sprawdzone. encrypt: k = secrets.randbelow(P - 2) + 1, zwraca pow(G, k, P) i M * pow(y, k, P) \% P. decrypt: s = pow(C1, x, P), potem C2 * pow(s, -1, P) \% P. Ostatnia linijka prowadzi do następnego slajdu.}
\end{frame}

% CHECKED (python3): M = 7, y = 19: k = 3 gives (10, 12), k = 7 gives (17, 13)
\begin{frame}{The same message twice}
  \begin{center}\small A new $k$ gives a \textbf{new ciphertext}, every time.\end{center}
  \begin{center}
    \begin{tikzpicture}[>=Stealth]
      \node[msgb, font=\small] (m) at (0,0) {$M = 7$};
      \node[sec, font=\scriptsize] (k1) at (2.6,0.75) {$k = 3$};
      \node[sec, font=\scriptsize] (k2) at (2.6,-0.75) {$k = 7$};
      \envpic{5.4}{0.75}{sBrown!15}{sBrown}
      \envpic{5.4}{-0.75}{sBrown!15}{sBrown}
      \node[font=\small\bfseries, text=nInk, fill=sBrown!15, inner sep=1pt] at (5.4,0.55) {(10, 12)};
      \node[font=\small\bfseries, text=nInk, fill=sBrown!15, inner sep=1pt] at (5.4,-0.95) {(17, 13)};
      \draw[data] (m.east) -- (k1.west); \draw[data] (m.east) -- (k2.west);
      \draw[data] (k1) -- (4.55,0.75); \draw[data] (k2) -- (4.55,-0.75);
    \end{tikzpicture}
  \end{center}
  \begin{center}
    \renewcommand{\arraystretch}{1.4}
    \begin{tabular}{lll}
      textbook RSA: & the same message, \textcolor{aWrite}{\textbf{the same ciphertext}} & \textcolor{gray}{(the PIN guessing of lecture 3)} \\
      ElGamal: & the same message, \textcolor{aExec}{\textbf{a different ciphertext}} & \\
    \end{tabular}
  \end{center}
  \begin{center}\small The price: the ciphertext is two numbers, \textbf{twice as long as $p$}.\end{center}
  \vfill
  \keymsg{ElGamal is randomised by design.}
  \say{W wykładzie 3 zgadywaliśmy PIN, bo podręcznikowe RSA zawsze daje ten sam szyfrogram. ElGamal tego problemu nie ma, bo losuje k.}
\end{frame}

% the first digits of the values in Task7/reuse.py
\begin{frame}{Task 7: k used twice}
  \vfill
  \taskgoal{recover the second message, without Bob's private key.}
  {\centering
    \begin{tikzpicture}
      \node[anchor=east, font=\small, text=nInk] at (-0.2,0.5) {first:};
      \node[ctx, font=\ttfamily\scriptsize, draw=AccentOrange, very thick, fill=AccentOrange!15] at (1.8,0.5) {C1 = 268389893958\dots};
      \node[ctx, font=\ttfamily\scriptsize] at (5.6,0.5) {C2 = 178206596531\dots};
      \node[anchor=east, font=\small, text=nInk] at (-0.2,-0.5) {second:};
      \node[ctx, font=\ttfamily\scriptsize, draw=AccentOrange, very thick, fill=AccentOrange!15] at (1.8,-0.5) {C1 = 268389893958\dots};
      \node[ctx, font=\ttfamily\scriptsize] at (5.6,-0.5) {C2 = 190653831068\dots};
      \node[font=\scriptsize\bfseries, text=AccentOrange] at (1.8,-1.2) {the same $k$, the same $C_1$};
    \end{tikzpicture}\par}\vspace{0.1cm}
  \taskready{the two ciphertexts, and the first message (from a public receipt), in \texttt{reuse.py}.}
  \vfill
  \opentask{Task7}{about 7 minutes}
  \vfill
  \say{Około 7 minut. Nie podpowiadać od razu. Jeśli trzeba: co w obu C2 jest takie samo? Zamiana tekstu na liczbę jest gotowa w pliku.}
\end{frame}

% CHECKED: Task7/solution/reuse.py prints the line below
\samedot
\begin{frame}{Task 7: what we saw}
  \vfill
  \begin{center}
    \begin{tikzpicture}
      \node[draw=aWrite, very thick, fill=aWrite!8, rounded corners=3pt, inner sep=6pt, font=\ttfamily\small, text=nInk]
        {\textcolor{aWrite}{\faExclamationTriangle}\ second message: New door code for the server room: 4096.};
    \end{tikzpicture}\\[0.4cm]
    Both $C_2$ hide their message with \textbf{the same mask} $y^{k}$.\\[2pt]
    The known first message gives the mask. The mask gives the second message.
  \end{center}
  \vfill
  \curious{diffie-hellman/\#6}
  \keymsg{Bob's private key was never needed.}
  \say{Sprawdzone. mask = C2\_first * pow(M\_first, -1, P) \% P, potem M\_second = C2\_second * pow(mask, -1, P) \% P. Klucz prywatny Boba nie był potrzebny ani przez chwilę. W symulatorze z linku ten sam atak na małych liczbach i na liczbach z zadania.}
\end{frame}

\samedot
\begin{frame}{The rule: a new k for every message}
  \vfill
  \begin{center}
    \begin{tikzpicture}
      \node[text=sAmber, font=\fontsize{56}{56}\selectfont] at (0,0) {\faDiceFive};
      \node[font=\small\bfseries, text=sAmberText] at (0,-1.2) {new every time};
    \end{tikzpicture}
  \end{center}
  \vfill
  \begin{center}
    \large $k$ must be \textbf{random}, \textbf{secret} and \textbf{new for every message}.\\[4pt]
    \small Secret: if $k$ leaks, anyone computes $y^{k}$ and divides it out.\\[4pt]
    \small The same lesson as the keystream and the nonce in the stream cipher lectures.
  \end{center}
  \vfill
  \keymsg{A repeated k gives away the message, just like a repeated keystream.}
  \say{Keystream i nonce to pojęcia z wykładów o szyfrach strumieniowych: tam też powtórka jednorazowej wartości otwierała wiadomość.}
\end{frame}

% CHECKED (python3): 2 * 12 % 23 = 1; pow(pow(10, 15, 23), -1, 23) = 14; 1 * 14 % 23 = 14 = 2 * 7
\begin{frame}{Mallory doubles the message}
  \vfill
  \begin{center}
    \begin{tikzpicture}[>=Stealth]
      \envpic{0}{0}{sBrown!15}{sBrown}
      \node[font=\small\bfseries, text=nInk, fill=sBrown!15, inner sep=1pt] at (0,-0.22) {(10, 12)};
      \node[capt, below=0.6cm] at (0,0) {from Alice, $M = 7$};
      \mallory{3.4}{0}
      \node[draw=sPink, thick, fill=white, rounded corners=3pt, font=\small, text=nInk, anchor=south] at (3.4,0.75) {$C_2 \times 2$: $12 \times 2 \bmod 23 = 1$};
      \draw[data] (0.9,0) -- (2.7,0);
      \draw[data] (4.1,0) -- (5.9,0);
      \envpic{6.8}{0}{sBrown!15}{sBrown}
      \node[font=\small\bfseries, text=nInk, fill=sBrown!15, inner sep=1pt] at (6.8,-0.22) {(10, \textcolor{sPinkText}{1})};
      \draw[data] (7.7,0) -- (9.2,0);
      \bob{10.0}{0}
      \node[msgb, font=\small, anchor=north] at (10.0,-0.95) {$M = 1 \times 14 \bmod 23 = \mathbf{14}$};
    \end{tikzpicture}
  \end{center}
  \vfill
  \begin{center}
    $C_2 = M \times y^{k}$, so $2 \times C_2 = (2M) \times y^{k}$: a valid ciphertext of $2M$.\\[2pt]
    Mallory knows neither $M$ nor any key, yet Bob decrypts $14 = 2 \times 7$.\\[4pt]
    \small A payment of 250 EUR could arrive as 500 EUR.
  \end{center}
  \vfill
  \keymsg{Textbook ElGamal hides $M$, but anyone can multiply it.}
  \say{Mallory nie odszyfrowuje niczego, tylko mnoży C2 przez 2. Bob liczy jak zwykle: s = 10\textasciicircum 15 mod 23 = 5, odwrotność 14, M = 1 · 14 = 14. Taka własność to podatność na modyfikację (malleability); podręcznikowe RSA ma ją tak samo (szyfrogram razy 2\textasciicircum e daje 2M). Dlatego ElGamala w praktyce nie używa się tak ,,goło'': wspólny sekret idzie do KDF, a wiadomość szyfruje szyfr z uwierzytelnieniem (DHIES, ECIES, HPKE z RFC 9180); o KDF w następnym wykładzie. Na tę zmianę nie pomaga też losowe k.}
\end{frame}

\begin{frame}{ElGamal and RSA side by side}
  \vfill
  \begin{center}
    \renewcommand{\arraystretch}{1.45}
    \begin{tabular}{lll}
      \toprule
      & \textbf{ElGamal} & \textbf{RSA} \\
      \midrule
      hard problem & discrete logarithm & factoring \\
      ciphertext & two numbers $(C_1, C_2)$ & one number \\
      randomness & built in (a new $k$) & needs padding (OAEP) \\
      \textcolor{aWrite}{\faExclamationTriangle}\ danger & \textcolor{aWrite}{a repeated $k$} & \textcolor{aWrite}{textbook use without padding} \\
      \textcolor{aWrite}{\faExclamationTriangle}\ textbook form & \textcolor{aWrite}{anyone can multiply $M$} & \textcolor{aWrite}{anyone can multiply $M$} \\
      in GnuPG (e-mail) & default encryption key until 2009 & default since 2009 \\
      \bottomrule
    \end{tabular}
  \end{center}
  \vfill
  \keymsg{Two ways to public-key crypto, each with its own trap.}
  % CHECKED: GnuPG NEWS: 0.3.0 "Key generation defaults to DSA/ElGamal keys"; 1.4.10 (2009-09-02) and 2.0.13 (2009-09-04):
  %   "2048 bit RSA keys are now generated by default" (github.com/gpg/gnupg, NEWS of STABLE-BRANCH-1-4 and 2-0);
  %   2.3.0 (2021-04-07): "gpg: Switch to ed25519/cv25519 as default public key algorithms"
  %   (lists.gnupg.org/pipermail/gnupg-announce/2021q2/000458.html)
  \say{Faktoryzacja to rozkład liczby na czynniki pierwsze, problem z wykładu 3. OAEP to losowe dopełnienie dla RSA, też z wykładu 3. GnuPG to najpopularniejszy program szyfrowanej poczty (standard OpenPGP): od pierwszych wersji (koniec lat 90.) domyślnie tworzył klucz ElGamal do szyfrowania, od wersji 1.4.10 i 2.0.13 (wrzesień 2009) domyślnie RSA, a od 2.3 (2021) krzywe eliptyczne, o których następny wykład.}
\end{frame}

\end{document}
